\documentclass[11pt]{article}
\usepackage[margin=1in]{geometry}
\usepackage{booktabs}
\usepackage{longtable}
\usepackage{array}
% ragged-right and right-aligned paragraph columns for the generated tables
\newcolumntype{L}[1]{>{\raggedright\arraybackslash}p{#1}}
\newcolumntype{R}[1]{>{\raggedleft\arraybackslash}p{#1}}
\usepackage{amsmath}
\usepackage{graphicx}
\usepackage{hyperref}
\hypersetup{colorlinks=true, linkcolor=blue, urlcolor=blue}
\sloppy

% Every number in this document is defined by analyze_datasf_catalogue.py in the two
% generated files below. The relevance flags are judgement, written by hand in
% relevance.csv beside this file, and are labelled as judgement wherever they appear.
% GENERATED BY analyze_datasf_catalogue.py --- do not edit by hand.
\newcommand{\catApiTotal}{273}
\newcommand{\catCached}{36}
\newcommand{\catCharts}{1}
\newcommand{\catCore}{35}
\newcommand{\catDatasets}{155}
\newcommand{\catDepartments}{23}
\newcommand{\catDomainsOnly}{2}
\newcommand{\catDomainsOnlyOld}{2}
\newcommand{\catFiles}{2}
\newcommand{\catFilters}{22}
\newcommand{\catGeom}{130}
\newcommand{\catKeyAny}{61}
\newcommand{\catKeyBlklot}{34}
\newcommand{\catKeyBlockLot}{30}
\newcommand{\catKeyCase}{18}
\newcommand{\catKeyPermit}{20}
\newcommand{\catLinks}{13}
\newcommand{\catMaps}{78}
\newcommand{\catNonDatasets}{118}
\newcommand{\catNotRelevant}{155}
\newcommand{\catOldHostCount}{273}
\newcommand{\catOldHostSameIds}{the same}
\newcommand{\catPlanningDept}{158}
\newcommand{\catPossible}{83}
\newcommand{\catPossibleNotCached}{82}
\newcommand{\catRetrieved}{2026-09-11}
\newcommand{\catRowsTotal}{28{,}161{,}918}
\newcommand{\catStale}{174}
\newcommand{\catStaleCutoff}{2024-09-11}
\newcommand{\catStaleYears}{2}
\newcommand{\catStories}{2}
\newcommand{\catTotal}{273}
\newcommand{\catUnassessed}{0}
\newcommand{\catWithRows}{255}
\newcommand{\catInvCommission}{1}
\newcommand{\catInvPipelineInternal}{4}
\newcommand{\catInvPlanning}{109}
\newcommand{\catInvPublished}{89}
\newcommand{\catInvUnpublished}{20}
\newcommand{\catRelPagingDistinct}{261}
\newcommand{\catRelPagingRows}{273}


\title{Everything DataSF Returns for ``Planning'':\\
\large A catalogue of the portal, and what in it the pipeline does not yet use}
\author{Dan Post}
\date{September 2026}

\begin{document}
\maketitle

\begin{abstract}
\noindent
The data-acquisition memo went to DataSF with a list of datasets it wanted. This memo goes
the other way: it asks the portal's catalogue for everything matching ``planning'' and
records what comes back, so that nothing useful is missed because nobody thought to ask for
it. Retrieved \catRetrieved.

\medskip\noindent
\textbf{The catalogue returns \catTotal\ assets, and the site's figure of 155 matches its
datasets alone.} The Discovery API counts \catDatasets\ datasets, \catMaps\ maps, \catFilters\
filtered views, \catLinks\ external links, \catStories\ stories, \catFiles\ files and
\catCharts\ chart; most of the \catNonDatasets\ that are not datasets are map views and
filtered views of datasets. Two facts about the
API are worth recording because both produce silently wrong answers: paged in its default
relevance order it returned \catRelPagingRows\ rows holding only \catRelPagingDistinct\
distinct assets, and the old host name \texttt{data.sfgov.org} is \emph{not} a dead end ---
passed as \texttt{search\_context} it returns \catOldHostCount\ results, \catOldHostSameIds\
assets as the new one.

\medskip\noindent
\textbf{There is still no dataset of Commission actions}, published or unpublished: not in
the \catTotal\ results, and not among the Planning Department's \catInvPlanning\ entries in
the city's own dataset inventory. The conditions memo's negative stands.

\medskip\noindent
\textbf{What the pipeline is not using.} \catPossibleNotCached\ assets are flagged, by
judgement, as core or possibly useful and are not cached. The two that matter most are
delay data nobody has measured against: the Planning Department's record-level review-time
metrics, keyed on the Accela record number, and DBI's station-by-station permit routing,
which dates every review step including Planning's. A newer publication of the current
zoning layer also exists beside the one the zoning panel carries forward.
\end{abstract}

\tableofcontents
\newpage

\section{The query, and why the count is not 155}

The brief records that the portal's own search reports 155 results for ``planning''. The
Socrata Discovery API, asked the same question with the domain set to \texttt{data.sf.gov}
and paged to exhaustion, returns \textbf{\catApiTotal}. Table~\ref{tab:types} reconciles the
two by asset type: the API's own facet count for datasets is \catDatasets, and when the full
result set is paged in a stable order the assets typed ``dataset'' are exactly those
\catDatasets. The site's figure matches that datasets count, and the other
\catNonDatasets\ results are \catMaps\ maps, \catFilters\ filtered views, \catLinks\
external links, \catStories\ stories, \catFiles\ files and \catCharts\ chart.

That reading is \emph{consistent with} the number the brief reports and has not been
verified against the page itself: the portal's search is a JavaScript application whose
count is filled in after the page loads, and no JavaScript engine was available to render
it. What is verified is the API's side of the equation, type by type.

% GENERATED BY analyze_datasf_catalogue.py --- do not edit by hand.
\begin{table}[htbp]\centering
\caption{The catalogue's answer to \texttt{q=planning}, by asset type. `Facet' is the API's own count with \texttt{only=} set to that type; `results' counts the assets returned when the full result set is paged in a stable order. The two agree type by type.}\label{tab:types}
\begin{tabular}{lrr}\toprule
Asset type & Facet count & Results\\\midrule
dataset & 155 & 155\\
map & 78 & 78\\
filtered view & 22 & 22\\
external link & 13 & 13\\
story & 2 & 2\\
file & 2 & 2\\
chart & 1 & 1\\
\midrule All & & 273\\
\bottomrule\end{tabular}\end{table}

\begin{table}[htbp]\centering
\caption{Which parameter names the host decides what the catalogue returns. Each row is the same query, \texttt{q=planning}, with the host passed differently.}\label{tab:hosttest}
\begin{tabular}{lr}\toprule
Query & Results\\\midrule
\texttt{domains + search\_context = data.sf.gov} & 273\\
\texttt{domains + search\_context = data.sfgov.org} & 273\\
\texttt{search\_context = data.sf.gov only} & 273\\
\texttt{search\_context = data.sfgov.org only} & 273\\
\texttt{domains = data.sf.gov only} & 2\\
\texttt{domains = data.sfgov.org only} & 2\\
\bottomrule\end{tabular}\end{table}

\begin{table}[htbp]\centering
\caption{By publishing department and by type. `Stale' is not updated in the 2 years before retrieval; `rows' sums \texttt{count(*)} over the tabular assets. Departments with fewer than three assets are pooled.}\label{tab:depts}
\resizebox{\textwidth}{!}{%
\begin{tabular}{lrrrrrr}\toprule
Department & Assets & Datasets & Maps & Other & Stale & Rows\\\midrule
Planning & 158 & 85 & 59 & 14 & 125 & 6{,}257{,}779\\
Municipal Transportation Agency & 15 & 9 & 4 & 2 & 12 & 22{,}693\\
Building Inspection & 14 & 1 & 2 & 11 & 0 & 13{,}511{,}840\\
Office of Resilience and Capital Planning & 14 & 12 & 1 & 1 & 13 & 4{,}061\\
(departments with fewer than three) & 12 & 9 & 2 & 1 & 6 & 4{,}248{,}182\\
Public Works & 11 & 9 & 2 & 0 & 0 & 1{,}450{,}981\\
GSA - City Administrator's Office & 8 & 5 & 2 & 1 & 6 & 239{,}131\\
Environment & 7 & 4 & 1 & 2 & 1 & 68{,}269\\
(not stated) & 6 & 2 & 3 & 1 & 5 & 4{,}330\\
Controller & 5 & 5 & 0 & 0 & 1 & 1{,}599{,}129\\
Technology & 5 & 3 & 1 & 1 & 2 & 474{,}211\\
Human Resources & 4 & 4 & 0 & 0 & 1 & 49{,}495\\
Ethics Commission & 4 & 3 & 0 & 1 & 0 & 2{,}049\\
Mayor & 4 & 3 & 1 & 0 & 0 & 2{,}043\\
Health Service System & 3 & 0 & 0 & 3 & 1 & 2{,}003\\
DataSF & 3 & 1 & 0 & 2 & 1 & 225{,}722\\
\bottomrule\end{tabular}}\end{table}

\begin{table}[htbp]\centering
\caption{By the catalogue's own category.}\label{tab:cats}
\begin{tabular}{lrr}\toprule
Category & Assets & of which datasets\\\midrule
(none) & 85 & 43\\
COVID-19 & 2 & 1\\
City Infrastructure & 21 & 15\\
City Management and Ethics & 14 & 12\\
Culture and Recreation & 6 & 3\\
Economy and Community & 11 & 4\\
Energy and Environment & 15 & 9\\
Geographic Locations and Boundaries & 46 & 27\\
Health and Social Services & 8 & 5\\
Housing and Buildings & 39 & 18\\
Public Safety & 10 & 9\\
Transportation & 16 & 9\\
\bottomrule\end{tabular}\end{table}

\begin{longtable}{@{}L{3.2cm}L{1.95cm}R{1.6cm}L{1.8cm}L{5.2cm}@{}}
\caption{Flagged \emph{core} or \emph{possibly useful} and not already cached. The flag and the reason are judgement.}\label{tab:possible}\\\toprule
Asset & Identifier & Rows & Keys (by name) & Why\\\midrule\endfirsthead
\toprule Asset & Identifier & Rows & Keys (by name) & Why\\\midrule\endhead
Building Permit Addenda with Routing \emph{(core)} & \texttt{87xy-gk8d} & 3{,}998{,}159 & permit, date & Station-by-station permit routing with dates, including Planning's review stations: the permit clock decomposed.\\
Planning Department Project Application Review metrics \emph{(core)} & \texttt{d4jk-jw33} & 67{,}265 & case/\allowbreak record, date & Record-level Planning review-time metrics keyed on the Accela record number: the delay the paper measures, from the department's own clock.\\
Building Permit Application Issuance Metrics \emph{(possibly useful)} & \texttt{gzxm-jz5j} & 148{,}951 & block+lot, permit, date & DBI's own filing-to-issuance and review-step metrics by permit; a check on the permit clock.\\
Building Permit Application Review Metrics \emph{(possibly useful)} & \texttt{5bat-azvb} & 453{,}703 & block+lot, permit, date & DBI's own filing-to-issuance and review-step metrics by permit; a check on the permit clock.\\
Building Permits Deduplicated on Primary Address \emph{(possibly useful)} & \texttt{f2jc-ivnc} & 1{,}149{,}358 & block+lot, case/\allowbreak record, permit, date & DBI's own deduplication on primary address; a check on the earliest-filed collapse rule.\\
Building Permits with Contractor, Subcontractor, Architect, Designer, or Engineer \emph{(possibly useful)} & \texttt{cw8k-gwb7} & 853{,}025 & permit, date & Permits with their architects, engineers and expediters: the repeat players in the permit market.\\
Building Permits with Permit Contacts \emph{(possibly useful)} & \texttt{9itm-3rmi} & 1{,}725{,}431 & block+lot, case/\allowbreak record, permit, date & Permits with their architects, engineers and expediters: the repeat players in the permit market.\\
Dataset inventory \emph{(possibly useful)} & \texttt{y8fp-fbf5} & 1{,}336 & date & The city's inventory of candidate datasets, published or not: where an unpublished Commission-actions table would appear.\\
Datasets in 'Scoping' or 'In Progress' phase \emph{(possibly useful)} & \texttt{xu9j-dnfu} & 64 & date & The city's inventory of candidate datasets, published or not: where an unpublished Commission-actions table would appear.\\
Land Use with Active Parcel Information \emph{(possibly useful)} & \texttt{h5vv-2u24} & 224{,}322 & blklot & The 2023 land use pre-joined to active parcels.\\
SFEC Form 3410A - Registration Report for Permit Consultants \emph{(possibly useful)} & \texttt{umwe-sn9p} & 170 & date & Registered permit consultants and their filings: who is paid to move permits.\\
SFEC Form 3410B - Quarterly Report for Permit Consultants Filings \emph{(possibly useful)} & \texttt{k2rm-926a} & 1{,}531 & date & Registered permit consultants and their filings: who is paid to move permits.\\
SFEC Form 3500 - Disclosure Report for Developers \emph{(possibly useful)} & \texttt{bk74-rima} & 174 & case/\allowbreak record, date & Disclosures by developers of major projects, keyed on the planning case number: the political economy of large entitlements.\\
SFEC Form 3500 - Disclosure Report for Developers of Major City Projects Filings \emph{(possibly useful)} & \texttt{8je3-gv3q} & 174 & case/\allowbreak record, date & Disclosures by developers of major projects, keyed on the planning case number: the political economy of large entitlements.\\
108" Inundation Vulnerability Zone Line (Sea Level Rise + 100YR Flood Event) \emph{(possibly useful)} & \texttt{92e4-7ptg} & 1{,}005 &  & A sea-level-rise or climate overlay: a long-run cost shifter; check against demand\_estimation first.\\
Parcels with overlay attributes \emph{(possibly useful)} & \texttt{9grn-xjpx} & 236{,}556 & blklot, block+lot, date & Every parcel with its current overlays pre-joined; a cross-check on the spatial join.\\
Mayor's Office of Housing and Community Development Affordable Housing Portfolio \emph{(possibly useful)} & \texttt{pyxv-n29e} & 850 & date & How projects met the inclusionary requirement, and the affordable portfolio, with case or permit numbers.\\
Mayor's Office of Housing and Community Development Non-Unit Inclusionary Housing Alternatives \emph{(possibly useful)} & \texttt{g6kh-9pnv} & 153 & case/\allowbreak record, permit, date & How projects met the inclusionary requirement, and the affordable portfolio, with case or permit numbers.\\
100-Year Storm + 24" Sea Level Rise \emph{(possibly useful)} & \texttt{esku-ejgv} & 7 &  & A physical-hazard zone: a construction-cost shifter; check against demand\_estimation's hazard layers first.\\
100-Year Storm Flood Risk Zone (July 2022) \emph{(possibly useful)} & \texttt{jzu3-4yxp} & 1 & date & A physical-hazard zone: a construction-cost shifter; check against demand\_estimation's hazard layers first.\\
100-year storm + 66 inches Sea Level Rise \emph{(possibly useful)} & \texttt{6ann-usi8} & 1 &  & A physical-hazard zone: a construction-cost shifter; check against demand\_estimation's hazard layers first.\\
Dam or Reservoir Failure Hazard Zone \emph{(possibly useful)} & \texttt{7crf-46s6} & 7 & date & A physical-hazard zone: a construction-cost shifter; check against demand\_estimation's hazard layers first.\\
FEMA FIRM Flood Hazards (Coastal) - 2021 Update \emph{(possibly useful)} & \texttt{jyce-e25k} & 222 &  & A physical-hazard zone: a construction-cost shifter; check against demand\_estimation's hazard layers first.\\
HAQR Priority Green Infrastructure Zones \emph{(possibly useful)} & \texttt{nn26-kuy2} & 1 &  & A sea-level-rise or climate overlay: a long-run cost shifter; check against demand\_estimation first.\\
Hayward M7.\allowbreak{}0 Earthquake Scenario \emph{(possibly useful)} & \texttt{d6pi-gknj} & 26 & date & A physical-hazard zone: a construction-cost shifter; check against demand\_estimation's hazard layers first.\\
Landslide Susceptibility Hazard Zones \emph{(possibly useful)} & \texttt{bna4-itif} & 8 & date & A physical-hazard zone: a construction-cost shifter; check against demand\_estimation's hazard layers first.\\
San Andreas M7.\allowbreak{}8 Earthquake Scenario \emph{(possibly useful)} & \texttt{akmp-ztcg} & 27 & date & A physical-hazard zone: a construction-cost shifter; check against demand\_estimation's hazard layers first.\\
Soil Liquefaction Hazard Zone \emph{(possibly useful)} & \texttt{i4t7-35u3} & 2 &  & A physical-hazard zone: a construction-cost shifter; check against demand\_estimation's hazard layers first.\\
Tsunami Inundation Hazard Zone (2021 Update) \emph{(possibly useful)} & \texttt{7p2k-c3z6} & 1 &  & A physical-hazard zone: a construction-cost shifter; check against demand\_estimation's hazard layers first.\\
Building Permit Completeness Check Review Metrics \emph{(possibly useful)} & \texttt{abh5-gwaq} & 6{,}618 & date & DBI's own filing-to-issuance and review-step metrics by permit; a check on the permit clock.\\
Permit\allowbreak{}SF Permitting Data \emph{(possibly useful)} & \texttt{tyz3-vt28} & 4{,}999 & blklot, case/\allowbreak record, date & PermitSF submission-to-issuance times with record numbers and parcels; recent and small.\\
AHBP\_\allowbreak{}eligible\_\allowbreak{}parcels \emph{(possibly useful)} & \texttt{fizh-zaxt} & 30{,}431 & blklot & Parcels eligible for the Affordable Housing Bonus Program: a dated policy treatment by parcel.\\
AHBP\_\allowbreak{}zoning\_\allowbreak{}districts \emph{(possibly useful)} & \texttt{y5fw-t3z4} & 736 &  & The zoning districts behind the Affordable Housing Bonus Program eligibility.\\
Air Pollutant Exposure Zone \emph{(possibly useful)} & \texttt{t65d-x6p8} & 1 &  & An overlay that triggers extra review (air-quality exposure, steep slope).\\
Analysis Neighborhoods \emph{(possibly useful)} & \texttt{j2bu-swwd} & 41 &  & Analysis neighbourhoods and their tract crosswalk: the geography the permit comparison uses.\\
Analysis Neighborhoods - 2010 census tracts assigned to neighborhoods \emph{(possibly useful)} & \texttt{m46u-xzix} & 195 &  & Analysis neighbourhoods and their tract crosswalk: the geography the permit comparison uses.\\
Analysis Neighborhoods - 2020 census tracts assigned to neighborhoods \emph{(possibly useful)} & \texttt{sevw-6tgi} & 242 & date & Analysis neighbourhoods and their tract crosswalk: the geography the permit comparison uses.\\
Article 11 Rating \emph{(possibly useful)} & \texttt{ejht-vp5d} & 786 & case/\allowbreak record, date & A historic designation that brings a project before a commission; overlaps the parcel-keyed historic status already listed.\\
Article 11 Ratings \emph{(possibly useful)} & \texttt{6m3x-8fu4} & 786 & blklot & Historic status by parcel or district: a regulatory constraint that triggers review.\\
Cultural Districts \emph{(possibly useful)} & \texttt{5xmc-5bjj} & 10 & date & Cultural districts: a geography used in land-use controls since 2018.\\
Downtown Plan Monitoring Report 2012 \emph{(possibly useful)} & \texttt{7d7s-7cc6} & --- &  & Downtown Plan monitoring: office-space totals behind the office-allocation docket.\\
Downtown Plan Monitoring Report 2012 \emph{(possibly useful)} & \texttt{eh2b-k988} & --- &  & Downtown Plan monitoring: office-space totals behind the office-allocation docket.\\
Downtown Plan Monitoring Report 2013 \emph{(possibly useful)} & \texttt{fhn6-zx2j} & --- &  & Downtown Plan monitoring: office-space totals behind the office-allocation docket.\\
Existing Commercial Wireless Telecommunication Services Facilities in San Francisco \emph{(possibly useful)} & \texttt{aa26-h926} & 883 & block+lot & Wireless facilities: wireless conditional uses are a large share of the CU docket.\\
Historic Context Statements \emph{(possibly useful)} & \texttt{43um-8u7u} & 5{,}503 & blklot & Parcels covered by adopted historic context statements: a review trigger by parcel.\\
Historic Districts \emph{(possibly useful)} & \texttt{63x5-g3m4} & 204 & date & Historic status by parcel or district: a regulatory constraint that triggers review.\\
Historic Resource Status by Parcel \emph{(possibly useful)} & \texttt{3tsw-4idn} & 163{,}278 & blklot, date & Historic status by parcel or district: a regulatory constraint that triggers review.\\
Historic Special Use Districts - 2010 \emph{(possibly useful)} & \texttt{6a3d-758c} & 79 &  & A dated special-use district vintage; the zoning panel carries only the current SUD layer.\\
Historic Special Use Districts - 2011 \emph{(possibly useful)} & \texttt{8vnp-htwm} & 106 &  & A dated special-use district vintage; the zoning panel carries only the current SUD layer.\\
Historic Special Use Districts - 2012 \emph{(possibly useful)} & \texttt{ff5j-uv3a} & 148 &  & A dated special-use district vintage; the zoning panel carries only the current SUD layer.\\
Historic Special Use Districts - 2014 \emph{(possibly useful)} & \texttt{b6kc-439r} & 150 &  & A dated special-use district vintage; the zoning panel carries only the current SUD layer.\\
Historic Special Use Districts - 2015 \emph{(possibly useful)} & \texttt{y7mn-8sik} & 146 &  & A dated special-use district vintage; the zoning panel carries only the current SUD layer.\\
Historic Special Use Districts -2013 \emph{(possibly useful)} & \texttt{hbc8-ncn7} & 147 &  & A dated special-use district vintage; the zoning panel carries only the current SUD layer.\\
Historic Zoning Maps \emph{(possibly useful)} & \texttt{ekvb-kch8} & --- &  & Scanned zoning maps from 1921 onward: the only pre-1998 zoning on the portal, as images.\\
Housing Balance March 2016 \emph{(possibly useful)} & \texttt{8iri-b2sz} & 2{,}637 & blklot, date & The 2016 housing balance snapshot by parcel: units added and lost, once.\\
Housing Sustainability Districts \emph{(possibly useful)} & \texttt{7way-43xh} & 1{,}678 &  & Housing Sustainability Districts: by-right approval geography, the ministerial counterfactual.\\
Landmark Districts \emph{(possibly useful)} & \texttt{knm6-5ej6} & 20 & case/\allowbreak record, date & A historic designation that brings a project before a commission; overlaps the parcel-keyed historic status already listed.\\
Landmarks \emph{(possibly useful)} & \texttt{rzic-39gi} & 362 &  & A historic designation that brings a project before a commission; overlaps the parcel-keyed historic status already listed.\\
Landmarks Listed in Article 10 of the San Francisco Planning Code \emph{(possibly useful)} & \texttt{97yj-54sx} & 370 & blklot & Historic status by parcel or district: a regulatory constraint that triggers review.\\
Maher \emph{(possibly useful)} & \texttt{v5ic-hjgw} & 974 &  & Sites subject to the Maher ordinance (contamination): a construction-cost shifter by site.\\
National Register Historic Districts \emph{(possibly useful)} & \texttt{4yr8-u35c} & 60 & date & A historic designation that brings a project before a commission; overlaps the parcel-keyed historic status already listed.\\
Neighborhood\allowbreak{}Quadrants \emph{(possibly useful)} & \texttt{6dn4-si7z} & 4 &  & A Planning Department administrative geography (district, area, team boundary), useful only as a grouping.\\
Planning Areas \emph{(possibly useful)} & \texttt{bkzf-eahq} & 72 &  & A Planning Department administrative geography (district, area, team boundary), useful only as a grouping.\\
Planning Department's Current Planning Team Boundaries \emph{(possibly useful)} & \texttt{ujhp-tya6} & 7 &  & A Planning Department administrative geography (district, area, team boundary), useful only as a grouping.\\
Planning Districts \emph{(possibly useful)} & \texttt{kyw5-5shs} & 17 &  & A Planning Department administrative geography (district, area, team boundary), useful only as a grouping.\\
Planning Neighborhood Groups Map \emph{(possibly useful)} & \texttt{mw29-m2za} & 37 &  & Neighbourhood-group notification boundaries: the geography of who is noticed.\\
Priority Development Areas \emph{(possibly useful)} & \texttt{k5ee-yrip} & 16 &  & Priority Development Areas: a regional policy geography.\\
Privately Owned Public Open Spaces \emph{(possibly useful)} & \texttt{65ik-7wqd} & 81 & block+lot & An obligation commonly imposed as a condition (open space, public art): a cross-check on condition content.\\
Proposed Commercial Wireless Telecommunication Services Facilities in San Francisco \emph{(possibly useful)} & \texttt{ewnu-9y3h} & 117 & block+lot & Wireless facilities: wireless conditional uses are a large share of the CU docket.\\
Public Art (from 1\% Art Program) \emph{(possibly useful)} & \texttt{cf6e-9e4j} & 65 &  & An obligation commonly imposed as a condition (open space, public art): a cross-check on condition content.\\
SF Find Neighborhoods \emph{(possibly useful)} & \texttt{gfpk-269f} & 117 &  & A Planning Department administrative geography (district, area, team boundary), useful only as a grouping.\\
San Francisco Environmental Justice Communities Map \emph{(possibly useful)} & \texttt{y6ci-vpnb} & 2{,}702 &  & Environmental-justice communities: a policy geography used in review.\\
San Francisco Land Use 2023 (No Geom) \emph{(possibly useful)} & \texttt{qbf8-heuf} & 154{,}476 & blklot & An archived land-use vintage; with the current layer it makes a land-use panel.\\
San Francisco Open Space for Shadow Study Analysis \emph{(possibly useful)} & \texttt{xk8z-bcqz} & --- &  & Open-space boundaries used in shadow studies: shadow is a recurring ground for discretionary review.\\
San Francisco Property Information Map \emph{(possibly useful)} & \texttt{i8ew-h6z7} & --- &  & A link to the Property Information Map, the per-parcel case and permit history; not a dataset.\\
Slopes 20 Percent Or Greater \emph{(possibly useful)} & \texttt{s66h-h4vp} & 28{,}055 & date & An overlay that triggers extra review (air-quality exposure, steep slope).\\
Zoning Map - Special Use Districts \emph{(possibly useful)} & \texttt{5yf5-ms5f} & 145 &  & A newer publication of the special-use district layer than the cached rxfc-8aap.\\
Zoning Map - Zoning Districts \emph{(possibly useful)} & \texttt{3i4a-hu95} & 10{,}617 &  & A newer publication of the current zoning layer (updated at retrieval; the cached xzez-p3nc was last updated 2024-02-23) --- the panel's carried-forward snapshot should be checked against it.\\
{[}ARCHIVED{]} San Francisco Land Use - 2020 \emph{(possibly useful)} & \texttt{ygi5-84iq} & 155{,}395 & blklot, block+lot & An archived land-use vintage; with the current layer it makes a land-use panel.\\
{[}ARCHIVED{]} San Francisco Land Use - 2023 \emph{(possibly useful)} & \texttt{fdfd-xptc} & 154{,}476 & blklot & An archived land-use vintage; with the current layer it makes a land-use panel.\\
City Lands \emph{(possibly useful)} & \texttt{gtnh-hgvs} & 1{,}035 & date & City-owned land: a cross-check on the risk set's public-land restriction.\\
San Francisco ZIP Codes \emph{(possibly useful)} & \texttt{uq3t-6t53} & 32 &  & ZIP code areas: the parcel-to-ZIP step the Zillow indices need.\\
\bottomrule\end{longtable}

\begin{longtable}{@{}L{5.3cm}L{1.95cm}L{2.0cm}L{1.8cm}L{2.6cm}@{}}
\caption{Stale: not updated in the 2 years before retrieval, oldest first.}\label{tab:stale}\\\toprule
Asset & Identifier & Type & Last update & Department\\\midrule\endfirsthead
\toprule Asset & Identifier & Type & Last update & Department\\\midrule\endhead
San Francisco Property Information Map & \texttt{i8ew-h6z7} & external link & 2012-09-26 & Planning\\
Commerce \& Industry Inventory 2011 Data & \texttt{cez4-442n} & external link & 2013-10-30 & Planning\\
Downtown Plan Monitoring Report 2012 & \texttt{eh2b-k988} & external link & 2013-11-13 & Planning\\
Commerce \& Industry Inventory 2012 Data & \texttt{uuiu-f3qh} & external link & 2013-11-14 & Planning\\
Downtown Plan Monitoring Report 2012 & \texttt{7d7s-7cc6} & external link & 2013-11-27 & Planning\\
Commerce \& Industry Inventory 2012 & \texttt{nrnx-fegt} & external link & 2013-11-27 & Planning\\
Commerce \& Industry Inventory 2011 & \texttt{ykqf-4jmt} & external link & 2013-11-27 & Planning\\
Downtown Plan Monitoring Report 2013 & \texttt{fhn6-zx2j} & external link & 2014-07-11 & Planning\\
2015 Project Appendix Map & \texttt{dyxq-hen7} & map & 2015-08-10 & GSA - City Administrator's Office\\
2015 Project Appendix & \texttt{xpdk-pzjd} & dataset & 2015-08-10 & GSA - City Administrator's Office\\
AHBP\_\allowbreak{}eligible\_\allowbreak{}parcels & \texttt{fizh-zaxt} & dataset & 2015-09-10 & Planning\\
Affordable Housing Bonus Program Zoning Districts and Eligible Parcels & \texttt{hxy5-cgnc} & map & 2015-09-10 & Planning\\
AHBP\_\allowbreak{}zoning\_\allowbreak{}districts & \texttt{y5fw-t3z4} & dataset & 2015-09-10 & Planning\\
Travel Decision Survey Data 2013 & \texttt{j6i5-prai} & file & 2015-09-25 & Municipal Transportation Agency\\
Neighborhood\allowbreak{}Quadrants & \texttt{6dn4-si7z} & dataset & 2016-03-06 & Planning\\
Planning Department Neighborhood Quadrants & \texttt{x89n-qawp} & map & 2016-03-06 & Planning\\
San Francisco Flood Health Vulnerability & \texttt{cne3-h93g} & dataset & 2016-03-25 & Public Health\\
Vehicle Miles Traveled Estimates from SF-CHAMP – 2012 & \texttt{g4ye-s75w} & external link & 2016-04-06 & Planning\\
Housing Balance March 2016 & \texttt{8iri-b2sz} & dataset & 2016-05-25 & Planning\\
Invest in Neighborhood Areas & \texttt{667d-aah6} & dataset & 2016-06-29 & Economic \& Workforce Development\\
Map of SF Urban Tree Canopy & \texttt{55pv-5zcc} & map & 2016-07-12 & Planning\\
Urban Tree Canopy & \texttt{ni2e-vpbg} & dataset & 2016-07-12 & Planning\\
Map of Historic Zoning Districts - 2008 & \texttt{ahwu-55yv} & map & 2016-08-04 & Planning\\
Green\_\allowbreak{}Roofs & \texttt{sfnk-6tdn} & dataset & 2016-08-04 & Planning\\
Historic\allowbreak{}Zoning\allowbreak{}Districts\allowbreak{}2008 & \texttt{tsdh-z53a} & dataset & 2016-08-04 & Planning\\
Civic Center Cultural Landscape Inventory & \texttt{r6rs-2mxu} & external link & 2016-09-27 & Planning\\
108" Inundation Vulnerability Zone Line (Sea Level Rise + 100YR Flood Event) & \texttt{92e4-7ptg} & dataset & 2017-03-20 & GSA - City Administrator's Office\\
Map of 108" Inundation Vulnerability Zone Line (Sea Level Rise + 100YR Flood Event) & \texttt{ff7h-99em} & map & 2017-03-20 & GSA - City Administrator's Office\\
Map of slopes of 20\% or greater & \texttt{3vv2-nvev} & map & 2017-08-10 & Planning\\
Slopes 20 Percent Or Greater & \texttt{s66h-h4vp} & dataset & 2017-08-10 & Planning\\
MTA.\allowbreak{}equitystrategy\_\allowbreak{}nhoods\_\allowbreak{}data & \texttt{qqn9-j3rz} & dataset & 2018-10-01 & Municipal Transportation Agency\\
MTA.\allowbreak{}paintedsafetyzones & \texttt{qqnk-bwbc} & map & 2018-10-01 & Municipal Transportation Agency\\
San Francisco Open Space for Shadow Study Analysis & \texttt{xk8z-bcqz} & external link & 2019-04-09 & Planning\\
Tall Building Inventory & \texttt{5kya-mfst} & dataset & 2019-07-11 & GSA - City Administrator's Office\\
Child Abuse Prevention Services in San Francisco & \texttt{3had-h899} & dataset & 2019-12-06 & Controller\\
San Andreas M7.\allowbreak{}8 Earthquake Scenario & \texttt{akmp-ztcg} & dataset & 2020-02-24 & Office of Resilience and Capital Planning\\
Hayward M7.\allowbreak{}0 Earthquake Scenario & \texttt{d6pi-gknj} & dataset & 2020-02-24 & Office of Resilience and Capital Planning\\
Soil Liquefaction Hazard Zone & \texttt{i4t7-35u3} & dataset & 2020-02-24 & Office of Resilience and Capital Planning\\
100-year storm + 66 inches Sea Level Rise & \texttt{6ann-usi8} & dataset & 2020-02-25 & Office of Resilience and Capital Planning\\
Dam or Reservoir Failure Hazard Zone & \texttt{7crf-46s6} & dataset & 2020-02-25 & Office of Resilience and Capital Planning\\
Landslide Susceptibility Hazard Zones & \texttt{bna4-itif} & dataset & 2020-02-25 & Office of Resilience and Capital Planning\\
100-Year Storm + 24" Sea Level Rise & \texttt{esku-ejgv} & dataset & 2020-02-25 & Office of Resilience and Capital Planning\\
Tsunami Inundation Hazard Zone (Outdated, 2015) & \texttt{q7uq-hghg} & dataset & 2020-02-25 & Office of Resilience and Capital Planning\\
City Construction Projects to Proceed During Shelter-in-place & \texttt{x83f-ux29} & dataset & 2020-04-10 & GSA - City Administrator's Office\\
Tenderloin COVID-19 Response: Sanitation Map & \texttt{fi63-bg5f} & map & 2020-06-23 & Emergency Management\\
COVID-19 Tenderloin Plan Zones & \texttt{xw26-dnmx} & dataset & 2020-06-23 & Emergency Management\\
FEMA FIRM Flood Hazards (Coastal) - 2021 Update & \texttt{jyce-e25k} & dataset & 2021-01-12 & Office of Resilience and Capital Planning\\
Priority Development Areas & \texttt{k5ee-yrip} & dataset & 2021-03-06 & Planning\\
Historic Zoning Maps & \texttt{ekvb-kch8} & external link & 2021-03-17 & Planning\\
Tsunami Inundation Hazard Zone (2021 Update) & \texttt{7p2k-c3z6} & dataset & 2021-08-03 & Office of Resilience and Capital Planning\\
Map of Tsunami Inundation Hazard Zone (2021 Update) & \texttt{xidx-bh43} & map & 2021-08-03 & Office of Resilience and Capital Planning\\
Existing Commercial Wireless Telecommunication Services Facilities in San Francisco & \texttt{aa26-h926} & dataset & 2021-09-17 & Planning\\
Proposed Commercial Wireless Telecommunication Services Facilities in San Francisco & \texttt{ewnu-9y3h} & dataset & 2021-09-17 & Planning\\
Travel Decision Survey 2021 & \texttt{b7eg-vqw2} & dataset & 2022-04-18 & Municipal Transportation Agency\\
SFMTA General Transit Feed Specification (GTFS) Production & \texttt{dni7-qpv3} & file & 2022-06-01 & Municipal Transportation Agency\\
Map of 2020 Census Tracts Assigned to Analysis Neighborhoods & \texttt{rqw6-h7c5} & map & 2022-07-08 & Planning\\
Analysis Neighborhoods - 2020 census tracts assigned to neighborhoods & \texttt{sevw-6tgi} & dataset & 2022-07-08 & Planning\\
100-Year Storm Flood Risk Zone (July 2022) & \texttt{jzu3-4yxp} & dataset & 2022-08-12 & Office of Resilience and Capital Planning\\
Planning Department's Current Planning Team Boundaries & \texttt{ujhp-tya6} & dataset & 2022-11-01 & Planning\\
City Employment Classification Definitions & \texttt{7br7-cj6v} & dataset & 2022-12-15 & Human Resources\\
San Francisco Environmental Justice Communities Map & \texttt{v25p-yprq} & map & 2023-07-21 & Planning\\
San Francisco Environmental Justice Communities Map & \texttt{y6ci-vpnb} & dataset & 2023-07-21 & Planning\\
ARCHIVED: Tenderloin Tent Count & \texttt{ubjr-y9qw} & dataset & 2023-07-24 & Other\\
San Francisco Aerial Photos (1938-2022) & \texttt{3keh-kpm9} & external link & 2023-10-02 & Planning\\
Historic Zoning Districts - March 2005 & \texttt{5uzm-v8n8} & dataset & 2023-10-16 & Planning\\
Historic Height \& Bulk Districts - 2014 & \texttt{6h9b-eksg} & map & 2023-10-16 & Planning\\
Historic Special Use Districts - 2011 & \texttt{8vnp-htwm} & dataset & 2023-10-16 & Planning\\
Historic Special Use Districts - 2012 & \texttt{9d9m-eyy3} & map & 2023-10-16 & Planning\\
Historic Height \& Bulk Districts - 2013 & \texttt{9mmf-fv7s} & dataset & 2023-10-16 & Planning\\
{[}DEPRECATED{]} Coastal Zone Area & \texttt{aeam-6mb5} & dataset & 2023-10-16 & Planning\\
Historic Zoning Districts - 1998 & \texttt{aypm-4d84} & map & 2023-10-16 & Planning\\
Historic Height \& Bulk Districts - 2013 & \texttt{bnc6-9btz} & map & 2023-10-16 & Planning\\
Map of Historic Zoning Districts - 2011 & \texttt{by3b-5pje} & map & 2023-10-16 & Planning\\
Historical Zoning Districts - April 2000 & \texttt{cfhj-8c6b} & dataset & 2023-10-16 & Planning\\
Historic Zoning Districts - March 2005 & \texttt{d7vm-pqzv} & map & 2023-10-16 & Planning\\
Areas Damaged by Fire Following 1906 Earthquake & \texttt{ff3a-iqhv} & dataset & 2023-10-16 & Planning\\
Historic Special Use Districts - 2012 & \texttt{ff5j-uv3a} & dataset & 2023-10-16 & Planning\\
Historic Height \& Bulk Districts - 2014 & \texttt{gu4h-44qp} & dataset & 2023-10-16 & Planning\\
Historic Special Use Districts -2013 & \texttt{hbc8-ncn7} & dataset & 2023-10-16 & Planning\\
Landmark Districts & \texttt{knm6-5ej6} & dataset & 2023-10-16 & Planning\\
Planning Districts & \texttt{kyw5-5shs} & dataset & 2023-10-16 & Planning\\
Historic Zoning Districts - 2007 & \texttt{pe87-v8tx} & dataset & 2023-10-16 & Planning\\
Historic Zoning Districts - 1998 & \texttt{piqy-pimd} & dataset & 2023-10-16 & Planning\\
Historical Zoning Districts - April 2000 & \texttt{prs8-k8k3} & map & 2023-10-16 & Planning\\
Green Connections Network & \texttt{q8g4-wqta} & map & 2023-10-16 & Planning\\
Map of Historic Special Use Districts -2013 & \texttt{r6jz-e8rw} & map & 2023-10-16 & Planning\\
Historic Zoning Districts - 2011 & \texttt{rt4q-mf68} & dataset & 2023-10-16 & Planning\\
Historic Zoning Districts - 2007 & \texttt{rxim-5dpc} & map & 2023-10-16 & Planning\\
Historic Special Use Districts - 2011 & \texttt{st6c-ij8w} & map & 2023-10-16 & Planning\\
Green Connections Network & \texttt{ta46-stby} & dataset & 2023-10-16 & Planning\\
Historic Height \& Bulk Districts - 2010 & \texttt{ti28-5szz} & map & 2023-10-16 & Planning\\
Planning Districts & \texttt{ttns-6zj3} & map & 2023-10-16 & Planning\\
Landmark Districts & \texttt{vnrd-fpg7} & map & 2023-10-16 & Planning\\
Historic Height \& Bulk Districts - 2010 & \texttt{xzb6-i4dc} & dataset & 2023-10-16 & Planning\\
Areas Damaged by Fire Following 1906 Earthquake & \texttt{yk2r-b4e8} & map & 2023-10-16 & Planning\\
Historic Zoning Districts - 2002 & \texttt{3yi7-eyfd} & dataset & 2023-10-17 & Planning\\
Historic Special Use Districts - 2010 & \texttt{6a3d-758c} & dataset & 2023-10-17 & Planning\\
Historic Zoning Districts - 2013 & \texttt{6jb9-g73z} & dataset & 2023-10-17 & Planning\\
Historic Zoning Districts - 2003 & \texttt{6ru7-zq4y} & dataset & 2023-10-17 & Planning\\
Historic Zoning Districts - 2001 & \texttt{8m6w-8yuk} & dataset & 2023-10-17 & Planning\\
Historic Zoning Districts - 2012 & \texttt{9xqx-n98b} & map & 2023-10-17 & Planning\\
Historic Height \& Bulk Districts - 2009 & \texttt{a4wu-zqx3} & dataset & 2023-10-17 & Planning\\
Historic Zoning Districts - Jan 2015 & \texttt{ada2-cu6t} & dataset & 2023-10-17 & Planning\\
Historic Zoning Districts - 2006 & \texttt{afkh-hfhr} & dataset & 2023-10-17 & Planning\\
Historic Zoning Districts - 2000 & \texttt{aksh-67x3} & map & 2023-10-17 & Planning\\
Historic Zoning Districts - 2005 & \texttt{b52k-gy2v} & dataset & 2023-10-17 & Planning\\
Historic Special Use Districts - 2014 & \texttt{b6kc-439r} & dataset & 2023-10-17 & Planning\\
Green Benefit Districts & \texttt{bmjt-kdwv} & map & 2023-10-17 & Planning\\
Analysis Neighborhoods - 2010 census tracts assigned to neighborhoods - Map & \texttt{bwbp-wk3r} & map & 2023-10-17 & Planning\\
Historic Special Use Districts - 2015 & \texttt{et47-hd49} & map & 2023-10-17 & Planning\\
Historic Zoning Districts - 2004 & \texttt{f883-k722} & map & 2023-10-17 & Planning\\
Historic Zoning Districts - 2002 & \texttt{ftvx-vtyc} & map & 2023-10-17 & Planning\\
SF Find Neighborhoods & \texttt{gfpk-269f} & dataset & 2023-10-17 & Planning\\
Historic Zoning Districts - 2014 & \texttt{hg44-eza7} & dataset & 2023-10-17 & Planning\\
Green Benefit Districts & \texttt{hyj5-mm3b} & dataset & 2023-10-17 & Planning\\
Planning Neighborhood Groups Map & \texttt{iacs-ws63} & map & 2023-10-17 & Planning\\
Historic Zoning Districts - 2000 & \texttt{itx2-wzp5} & dataset & 2023-10-17 & Planning\\
Analysis Neighborhoods & \texttt{j2bu-swwd} & dataset & 2023-10-17 & Planning\\
Historic Zoning Districts - 2014 & \texttt{jn35-yfmd} & map & 2023-10-17 & Planning\\
Historic Zoning Districts - 2009 & \texttt{jud5-ja46} & dataset & 2023-10-17 & Planning\\
Historic Special Use Districts - 2010 & \texttt{kmx9-ph84} & map & 2023-10-17 & Planning\\
Historic Zoning Districts - 2003 & \texttt{kspe-fmej} & map & 2023-10-17 & Planning\\
Analysis Neighborhoods - 2010 census tracts assigned to neighborhoods & \texttt{m46u-xzix} & dataset & 2023-10-17 & Planning\\
Historic Zoning Districts - 2013 & \texttt{manb-7mxp} & map & 2023-10-17 & Planning\\
Planning Neighborhood Groups Map & \texttt{mw29-m2za} & dataset & 2023-10-17 & Planning\\
Historic Height \& Bulk Districts - 2012 & \texttt{nh4m-jbyj} & dataset & 2023-10-17 & Planning\\
Analysis Neighborhoods & \texttt{p5b7-5n3h} & map & 2023-10-17 & Planning\\
Historic Zoning Districts - 2001 & \texttt{pdvd-w2q4} & map & 2023-10-17 & Planning\\
SF Find Neighborhoods & \texttt{pty2-tcw4} & map & 2023-10-17 & Planning\\
Historic Height \& Bulk Districts - 2011 & \texttt{qcxd-tp4u} & map & 2023-10-17 & Planning\\
Historic Height \& Bulk Districts - 2011 & \texttt{rwdp-2k4t} & dataset & 2023-10-17 & Planning\\
Historic Zoning Districts - 2004 & \texttt{s3eh-jdwb} & dataset & 2023-10-17 & Planning\\
Historic Height \& Bulk Districts - 2009 & \texttt{thsi-uh4u} & map & 2023-10-17 & Planning\\
Historic Height \& Bulk Districts - 2012 & \texttt{u3m4-4qjy} & map & 2023-10-17 & Planning\\
Historic Zoning Districts - 2009 & \texttt{u4br-5hb8} & map & 2023-10-17 & Planning\\
Historic Zoning Districts - 2005 & \texttt{u8bf-s4hg} & map & 2023-10-17 & Planning\\
Map of Historic Special Use Districts - 2014 & \texttt{uihu-8bih} & map & 2023-10-17 & Planning\\
{[}ARCHIVED{]} Map of Land Use - 2020 & \texttt{us3s-fp9q} & map & 2023-10-17 & Planning\\
Historic Zoning Districts - Jan 2015 & \texttt{utj8-tgqr} & map & 2023-10-17 & Planning\\
Map of Historic Zoning Districts - 2006 & \texttt{v9w2-c8v7} & map & 2023-10-17 & Planning\\
Historic Zoning Districts - 2012 & \texttt{vfz8-awmy} & dataset & 2023-10-17 & Planning\\
Historic Zoning Districts - 2010 & \texttt{w3j2-4hed} & map & 2023-10-17 & Planning\\
Historic Zoning Districts - 2010 & \texttt{x4gj-zjx7} & dataset & 2023-10-17 & Planning\\
Historic Special Use Districts - 2015 & \texttt{y7mn-8sik} & dataset & 2023-10-17 & Planning\\
{[}ARCHIVED{]} San Francisco Land Use - 2020 & \texttt{ygi5-84iq} & dataset & 2023-10-17 & Planning\\
San Francisco ZIP Codes & \texttt{srq6-hmpi} & map & 2023-10-27 & Technology\\
San Francisco ZIP Codes & \texttt{uq3t-6t53} & dataset & 2023-10-27 & Technology\\
Retiree Pensions Annual Benefit Received & \texttt{84zg-myyk} & dataset & 2023-11-08 & Retirement System\\
National Register Historic Districts & \texttt{4yr8-u35c} & dataset & 2024-01-12 & Planning\\
Public Art (from 1\% Art Program) & \texttt{cf6e-9e4j} & dataset & 2024-01-13 & Planning\\
On-street Carshare Parking & \texttt{g2t6-cyw6} & dataset & 2024-01-18 & Municipal Transportation Agency\\
Slow Streets & \texttt{hkz3-itiu} & dataset & 2024-01-18 & Municipal Transportation Agency\\
Map of On-street Carshare Parking & \texttt{wr3x-ijzr} & map & 2024-01-18 & Municipal Transportation Agency\\
Article 11 Rating & \texttt{ee6v-jtbt} & map & 2024-02-21 & (not stated)\\
Article 11 Rating & \texttt{ejht-vp5d} & dataset & 2024-02-21 & Planning\\
Zoning Map - Zoning Districts & \texttt{ibu8-4ccn} & map & 2024-02-23 & (not stated)\\
Zoning Map - Zoning Districts & \texttt{xzez-p3nc} & dataset & 2024-02-23 & (not stated)\\
{[}ARCHIVED{]} San Francisco Land Use - 2023 & \texttt{fdfd-xptc} & dataset & 2024-03-07 & Planning\\
Land Use with Active Parcel Information & \texttt{h5vv-2u24} & filtered view & 2024-03-07 & DataSF\\
Map of San Francisco Land Use - 2023 & \texttt{k8rg-ihdq} & map & 2024-03-07 & Planning\\
San Francisco Land Use 2023 (No Geom) & \texttt{qbf8-heuf} & filtered view & 2024-03-07 & Planning\\
Zoning Map - Special Use Districts & \texttt{rxfc-8aap} & dataset & 2024-05-01 & (not stated)\\
Zoning Map - Special Use Districts & \texttt{ry69-hnut} & map & 2024-05-01 & (not stated)\\
Landmarks & \texttt{rzic-39gi} & dataset & 2024-05-03 & Planning\\
Zoning Map - Special Sign Districts - Scenic Streets & \texttt{4hx7-bwhg} & map & 2024-05-08 & Planning\\
Zoning Map - Special Sign Districts - Scenic Streets & \texttt{4n3c-huha} & dataset & 2024-05-08 & Planning\\
SF Planning Permitting Data {[}Deprecated{]} & \texttt{kncr-c6jw} & dataset & 2024-05-15 & Planning\\
Traffic Signals & \texttt{ybh5-27n2} & dataset & 2024-05-31 & Municipal Transportation Agency\\
View of Medical and Dental Enrollment by Union for City \& County of SF Employees(CCSF) & \texttt{42ay-783m} & filtered view & 2024-06-13 & Health Service System\\
Sidewalk Widths (2014) & \texttt{4g86-grxu} & dataset & 2024-06-17 & Municipal Transportation Agency\\
San Francisco Communitywide Greenhouse Gas Inventory & \texttt{btm4-e4ak} & dataset & 2024-06-20 & Environment\\
HAQR Priority Green Infrastructure Zones & \texttt{nn26-kuy2} & dataset & 2024-06-20 & Office of Resilience and Capital Planning\\
Equity Strategy Neighborhoods & \texttt{f8bh-m8ez} & dataset & 2024-07-30 & Municipal Transportation Agency\\
Equity Strategy Neighborhoods & \texttt{u7ag-wi48} & map & 2024-07-30 & Municipal Transportation Agency\\
\bottomrule\end{longtable}

{\small\setlength{\tabcolsep}{4pt}
\begin{longtable}{@{}L{3.1cm}L{1.75cm}L{1.3cm}R{1.45cm}R{0.7cm}L{1.15cm}L{1.25cm}L{1.75cm}L{0.55cm}@{}}
\caption{The full catalogue, grouped by department. `Upd.' is the last data update; `Cols' the column count; `Keys' the pipeline join keys a column name suggests; `Flag' is judgement (c core, p possibly useful, n not relevant); a dagger marks a source the pipeline already caches.}\label{tab:full}\\\toprule
Asset & Identifier & Type & Rows & Col. & Upd. & Geom. & Keys & Fl.\\\midrule\endfirsthead
\toprule Asset & Identifier & Type & Rows & Col. & Upd. & Geom. & Keys & Fl.\\\midrule\endhead
\multicolumn{9}{@{}l}{\textbf{Planning} (158)}\\
AHBP\_\allowbreak{}eligible\_\allowbreak{}parcels & \texttt{fizh-zaxt} & dataset & 30{,}431 & 7 & 2015-09 & polygon & blklot & p\\
AHBP\_\allowbreak{}zoning\_\allowbreak{}districts & \texttt{y5fw-t3z4} & dataset & 736 & 8 & 2015-09 & polygon &  & p\\
Air Pollutant Exposure Zone & \texttt{t65d-x6p8} & dataset & 1 & 1 & 2025-03 & polygon &  & p\\
Analysis Neighborhoods & \texttt{j2bu-swwd} & dataset & 41 & 2 & 2023-10 & polygon &  & p\\
Analysis Neighborhoods - 2010 census tracts assigned to neighborhoods & \texttt{m46u-xzix} & dataset & 195 & 6 & 2023-10 & polygon &  & p\\
Analysis Neighborhoods - 2020 census tracts assigned to neighborhoods & \texttt{sevw-6tgi} & dataset & 242 & 12 & 2022-07 & polygon & date & p\\
Areas Damaged by Fire Following 1906 Earthquake & \texttt{ff3a-iqhv} & dataset & 1 & 9 & 2023-10 & polygon &  & n\\
Article 11 Rating & \texttt{ejht-vp5d} & dataset & 786 & 18 & 2024-02 & polygon & case/\allowbreak record, date & p\\
Article 11 Ratings & \texttt{6m3x-8fu4} & dataset & 786 & 8 & 2026-09 & polygon & blklot & p\\
Cultural Districts & \texttt{5xmc-5bjj} & dataset & 10 & 3 & 2026-09 & polygon & date & p\\
Existing Commercial Wireless Telecommunication Services Facilities in San Francisco & \texttt{aa26-h926} & dataset & 883 & 22 & 2021-09 &  & block+lot & p\\
Green Benefit Districts & \texttt{hyj5-mm3b} & dataset & 1 & 6 & 2023-10 & polygon &  & n\\
Green Connections Network & \texttt{ta46-stby} & dataset & 2{,}738 & 12 & 2023-10 & line &  & n\\
Green\_\allowbreak{}Roofs & \texttt{sfnk-6tdn} & dataset & 35 & 21 & 2016-08 & point &  & n\\
Historic Context Statements & \texttt{43um-8u7u} & dataset & 5{,}503 & 11 & 2026-09 & polygon & blklot & p\\
Historic Districts & \texttt{63x5-g3m4} & dataset & 204 & 13 & 2026-09 & polygon & date & p\\
Historic Height \& Bulk Districts - 2009$^{\dagger}$ & \texttt{a4wu-zqx3} & dataset & 6{,}703 & 6 & 2023-10 & polygon &  & c\\
Historic Height \& Bulk Districts - 2010$^{\dagger}$ & \texttt{xzb6-i4dc} & dataset & 7{,}130 & 6 & 2023-10 & polygon &  & c\\
Historic Height \& Bulk Districts - 2011$^{\dagger}$ & \texttt{rwdp-2k4t} & dataset & 7{,}145 & 6 & 2023-10 & polygon &  & c\\
Historic Height \& Bulk Districts - 2012$^{\dagger}$ & \texttt{nh4m-jbyj} & dataset & 7{,}195 & 6 & 2023-10 & polygon &  & c\\
Historic Height \& Bulk Districts - 2013$^{\dagger}$ & \texttt{9mmf-fv7s} & dataset & 7{,}314 & 6 & 2023-10 & polygon &  & c\\
Historic Height \& Bulk Districts - 2014$^{\dagger}$ & \texttt{gu4h-44qp} & dataset & 7{,}317 & 6 & 2023-10 & polygon &  & c\\
Historic Resource Status by Parcel & \texttt{3tsw-4idn} & dataset & 163{,}278 & 13 & 2026-09 & polygon & blklot, date & p\\
Historic Special Use Districts - 2010 & \texttt{6a3d-758c} & dataset & 79 & 5 & 2023-10 & polygon &  & p\\
Historic Special Use Districts - 2011 & \texttt{8vnp-htwm} & dataset & 106 & 6 & 2023-10 & polygon &  & p\\
Historic Special Use Districts - 2012 & \texttt{ff5j-uv3a} & dataset & 148 & 6 & 2023-10 & polygon &  & p\\
Historic Special Use Districts - 2014 & \texttt{b6kc-439r} & dataset & 150 & 6 & 2023-10 & polygon &  & p\\
Historic Special Use Districts - 2015 & \texttt{y7mn-8sik} & dataset & 146 & 6 & 2023-10 & polygon &  & p\\
Historic Special Use Districts -2013 & \texttt{hbc8-ncn7} & dataset & 147 & 6 & 2023-10 & polygon &  & p\\
Historic Zoning Districts - 1998$^{\dagger}$ & \texttt{piqy-pimd} & dataset & 152{,}569 & 35 & 2023-10 & polygon & blklot, block+lot & c\\
Historic Zoning Districts - 2000$^{\dagger}$ & \texttt{itx2-wzp5} & dataset & 153{,}292 & 32 & 2023-10 & polygon & blklot, block+lot & c\\
Historic Zoning Districts - 2001$^{\dagger}$ & \texttt{8m6w-8yuk} & dataset & 153{,}669 & 40 & 2023-10 & polygon & blklot, block+lot & c\\
Historic Zoning Districts - 2002$^{\dagger}$ & \texttt{3yi7-eyfd} & dataset & 153{,}663 & 40 & 2023-10 & polygon & blklot, block+lot & c\\
Historic Zoning Districts - 2003$^{\dagger}$ & \texttt{6ru7-zq4y} & dataset & 153{,}716 & 41 & 2023-10 & polygon & blklot, block+lot & c\\
Historic Zoning Districts - 2004$^{\dagger}$ & \texttt{s3eh-jdwb} & dataset & 191{,}062 & 48 & 2023-10 & polygon & blklot, block+lot, date & c\\
Historic Zoning Districts - 2005$^{\dagger}$ & \texttt{b52k-gy2v} & dataset & 196{,}344 & 26 & 2023-10 & polygon & blklot, block+lot & c\\
Historic Zoning Districts - 2006$^{\dagger}$ & \texttt{afkh-hfhr} & dataset & 23{,}132 & 6 & 2023-10 & polygon &  & c\\
Historic Zoning Districts - 2007$^{\dagger}$ & \texttt{pe87-v8tx} & dataset & 153{,}822 & 59 & 2023-10 & polygon & blklot, block+lot, date & c\\
Historic Zoning Districts - 2009$^{\dagger}$ & \texttt{jud5-ja46} & dataset & 23{,}355 & 7 & 2023-10 & polygon &  & c\\
Historic Zoning Districts - 2010$^{\dagger}$ & \texttt{x4gj-zjx7} & dataset & 9{,}897 & 7 & 2023-10 & polygon &  & c\\
Historic Zoning Districts - 2011$^{\dagger}$ & \texttt{rt4q-mf68} & dataset & 9{,}880 & 7 & 2023-10 & polygon &  & c\\
Historic Zoning Districts - 2012$^{\dagger}$ & \texttt{vfz8-awmy} & dataset & 9{,}889 & 7 & 2023-10 & polygon &  & c\\
Historic Zoning Districts - 2013$^{\dagger}$ & \texttt{6jb9-g73z} & dataset & 9{,}982 & 7 & 2023-10 & polygon &  & c\\
Historic Zoning Districts - 2014$^{\dagger}$ & \texttt{hg44-eza7} & dataset & 9{,}981 & 7 & 2023-10 & polygon &  & c\\
Historic Zoning Districts - Jan 2015$^{\dagger}$ & \texttt{ada2-cu6t} & dataset & 9{,}978 & 7 & 2023-10 & polygon &  & c\\
Historic Zoning Districts - March 2005 & \texttt{5uzm-v8n8} & dataset & 196{,}344 & 18 & 2023-10 & polygon & blklot, block+lot & n\\
Historic\allowbreak{}Zoning\allowbreak{}Districts\allowbreak{}2008$^{\dagger}$ & \texttt{tsdh-z53a} & dataset & 153{,}355 & 128 & 2016-08 & polygon & blklot, block+lot, date & c\\
Historical Zoning Districts - April 2000 & \texttt{cfhj-8c6b} & dataset & 153{,}292 & 13 & 2023-10 & polygon & blklot, block+lot & n\\
Housing Balance March 2016 & \texttt{8iri-b2sz} & dataset & 2{,}637 & 25 & 2016-05 & point & blklot, date & p\\
Housing Production - 2005-present$^{\dagger}$ & \texttt{xdht-4php} & dataset & 6{,}139 & 31 & 2026-09 &  & blklot, case/\allowbreak record, permit, date & c\\
Housing Sustainability Districts & \texttt{7way-43xh} & dataset & 1{,}678 & 3 & 2026-03 & polygon &  & p\\
Landmark Districts & \texttt{knm6-5ej6} & dataset & 20 & 16 & 2023-10 & polygon & case/\allowbreak record, date & p\\
Landmarks & \texttt{rzic-39gi} & dataset & 362 & 36 & 2024-05 & polygon &  & p\\
Landmarks Listed in Article 10 of the San Francisco Planning Code & \texttt{97yj-54sx} & dataset & 370 & 9 & 2026-09 & polygon & blklot & p\\
Maher & \texttt{v5ic-hjgw} & dataset & 974 & 4 & 2026-09 & polygon &  & p\\
National Register Historic Districts & \texttt{4yr8-u35c} & dataset & 60 & 34 & 2024-01 & polygon & date & p\\
Neighborhood-Specific Impact Fee Areas$^{\dagger}$ & \texttt{ntc3-dd64} & dataset & 33 & 7 & 2024-11 & polygon &  & c\\
Neighborhood\allowbreak{}Quadrants & \texttt{6dn4-si7z} & dataset & 4 & 8 & 2016-03 & polygon &  & p\\
Planning Areas & \texttt{bkzf-eahq} & dataset & 72 & 4 & 2026-09 & polygon &  & p\\
Planning Department Project Application Review metrics & \texttt{d4jk-jw33} & dataset & 67{,}265 & 14 & 2026-09 &  & case/\allowbreak record, date & c\\
Planning Department Records - Non-Projects$^{\dagger}$ & \texttt{y673-d69b} & dataset & 233{,}459 & 17 & 2026-09 &  & block+lot, case/\allowbreak record, permit, date & c\\
Planning Department Records - Projects$^{\dagger}$ & \texttt{qvu5-m3a2} & dataset & 54{,}522 & 107 & 2026-09 &  & block+lot, case/\allowbreak record, permit, date & c\\
Planning Department's Current Planning Team Boundaries & \texttt{ujhp-tya6} & dataset & 7 & 4 & 2022-11 & polygon &  & p\\
Planning Districts & \texttt{kyw5-5shs} & dataset & 17 & 5 & 2023-10 & polygon &  & p\\
Planning Neighborhood Groups Map & \texttt{mw29-m2za} & dataset & 37 & 2 & 2023-10 & polygon &  & p\\
Priority Development Areas & \texttt{k5ee-yrip} & dataset & 16 & 7 & 2021-03 & polygon &  & p\\
Privately Owned Public Open Spaces & \texttt{65ik-7wqd} & dataset & 81 & 31 & 2026-09 & point & block+lot & p\\
Proposed Commercial Wireless Telecommunication Services Facilities in San Francisco & \texttt{ewnu-9y3h} & dataset & 117 & 23 & 2021-09 &  & block+lot & p\\
Public Art (from 1\% Art Program) & \texttt{cf6e-9e4j} & dataset & 65 & 19 & 2024-01 & point &  & p\\
SF Find Neighborhoods & \texttt{gfpk-269f} & dataset & 117 & 3 & 2023-10 & polygon &  & p\\
SF Planning Permitting Data {[}Deprecated{]} & \texttt{kncr-c6jw} & dataset & 0 & 144 & 2024-05 & point & block+lot, case/\allowbreak record, permit, date & n\\
San Francisco Development Pipeline$^{\dagger}$ & \texttt{6jgi-cpb4} & dataset & 1{,}917 & 66 & 2026-08 & point & blklot, case/\allowbreak record, permit, date & c\\
San Francisco Environmental Justice Communities Map & \texttt{y6ci-vpnb} & dataset & 2{,}702 & 3 & 2023-07 & polygon &  & p\\
San Francisco Land Use$^{\dagger}$ & \texttt{c5ge-t6pj} & dataset & 153{,}765 & 23 & 2026-08 & polygon & blklot, date & p\\
Slopes 20 Percent Or Greater & \texttt{s66h-h4vp} & dataset & 28{,}055 & 7 & 2017-08 & polygon & date & p\\
Urban Tree Canopy & \texttt{ni2e-vpbg} & dataset & 289{,}219 & 4 & 2016-07 & polygon &  & n\\
Zoning Map - Height and Bulk Districts$^{\dagger}$ & \texttt{h9wh-cg3m} & dataset & 5{,}292 & 3 & 2026-09 & polygon &  & c\\
Zoning Map - Special Sign Districts & \texttt{g9f8-pkcz} & dataset & 13 & 4 & 2026-09 & polygon &  & n\\
Zoning Map - Special Sign Districts - Scenic Streets & \texttt{4n3c-huha} & dataset & 12 & 4 & 2024-05 & polygon &  & n\\
Zoning Map - Special Use Districts & \texttt{5yf5-ms5f} & dataset & 145 & 3 & 2026-09 & polygon &  & p\\
Zoning Map - Zoning Districts & \texttt{3i4a-hu95} & dataset & 10{,}617 & 7 & 2026-09 & polygon &  & p\\
{[}ARCHIVED{]} San Francisco Land Use - 2020 & \texttt{ygi5-84iq} & dataset & 155{,}395 & 27 & 2023-10 & polygon & blklot, block+lot & p\\
{[}ARCHIVED{]} San Francisco Land Use - 2023 & \texttt{fdfd-xptc} & dataset & 154{,}476 & 17 & 2024-03 & polygon & blklot & p\\
{[}DEPRECATED{]} Coastal Zone Area & \texttt{aeam-6mb5} & dataset & 1 & 5 & 2023-10 & polygon &  & n\\
{[}DEPRECATED{]} Historic Resources & \texttt{njrt-gtwr} & dataset & 0 & 15 & 2026-05 & polygon & blklot, date & n\\
Civic Center Cultural Landscape Inventory & \texttt{r6rs-2mxu} & external link & --- & 0 & 2016-09 &  &  & n\\
Commerce \& Industry Inventory 2011 & \texttt{ykqf-4jmt} & external link & --- & 0 & 2013-11 &  &  & n\\
Commerce \& Industry Inventory 2011 Data & \texttt{cez4-442n} & external link & --- & 0 & 2013-10 &  &  & n\\
Commerce \& Industry Inventory 2012 & \texttt{nrnx-fegt} & external link & --- & 0 & 2013-11 &  &  & n\\
Commerce \& Industry Inventory 2012 Data & \texttt{uuiu-f3qh} & external link & --- & 0 & 2013-11 &  &  & n\\
Downtown Plan Monitoring Report 2012 & \texttt{7d7s-7cc6} & external link & --- & 0 & 2013-11 &  &  & p\\
Downtown Plan Monitoring Report 2012 & \texttt{eh2b-k988} & external link & --- & 0 & 2013-11 &  &  & p\\
Downtown Plan Monitoring Report 2013 & \texttt{fhn6-zx2j} & external link & --- & 0 & 2014-07 &  &  & p\\
Historic Zoning Maps & \texttt{ekvb-kch8} & external link & --- & 0 & 2021-03 &  &  & p\\
San Francisco Aerial Photos (1938-2022) & \texttt{3keh-kpm9} & external link & --- & 0 & 2023-10 &  &  & n\\
San Francisco Open Space for Shadow Study Analysis & \texttt{xk8z-bcqz} & external link & --- & 0 & 2019-04 &  &  & p\\
San Francisco Property Information Map & \texttt{i8ew-h6z7} & external link & --- & 0 & 2012-09 &  &  & p\\
Vehicle Miles Traveled Estimates from SF-CHAMP – 2012 & \texttt{g4ye-s75w} & external link & --- & 0 & 2016-04 &  &  & n\\
San Francisco Land Use 2023 (No Geom) & \texttt{qbf8-heuf} & filtered view & 154{,}476 & 16 & 2024-03 &  & blklot & p\\
Affordable Housing Bonus Program Zoning Districts and Eligible Parcels & \texttt{hxy5-cgnc} & map & 736 & 0 & 2015-09 &  &  & n\\
Analysis Neighborhoods & \texttt{p5b7-5n3h} & map & 41 & 0 & 2023-10 &  &  & n\\
Analysis Neighborhoods - 2010 census tracts assigned to neighborhoods - Map & \texttt{bwbp-wk3r} & map & 195 & 0 & 2023-10 &  &  & n\\
Areas Damaged by Fire Following 1906 Earthquake & \texttt{yk2r-b4e8} & map & 1 & 0 & 2023-10 &  &  & n\\
Cultural Districts Map & \texttt{a5e9-skih} & map & 10 & 0 & 2026-09 &  &  & n\\
Green Benefit Districts & \texttt{bmjt-kdwv} & map & 1 & 0 & 2023-10 &  &  & n\\
Green Connections Network & \texttt{q8g4-wqta} & map & 2{,}738 & 0 & 2023-10 &  &  & n\\
Historic Height \& Bulk Districts - 2009 & \texttt{thsi-uh4u} & map & 6{,}703 & 0 & 2023-10 &  &  & n\\
Historic Height \& Bulk Districts - 2010 & \texttt{ti28-5szz} & map & 7{,}130 & 0 & 2023-10 &  &  & n\\
Historic Height \& Bulk Districts - 2011 & \texttt{qcxd-tp4u} & map & 7{,}145 & 0 & 2023-10 &  &  & n\\
Historic Height \& Bulk Districts - 2012 & \texttt{u3m4-4qjy} & map & 7{,}195 & 0 & 2023-10 &  &  & n\\
Historic Height \& Bulk Districts - 2013 & \texttt{bnc6-9btz} & map & 7{,}314 & 0 & 2023-10 &  &  & n\\
Historic Height \& Bulk Districts - 2014 & \texttt{6h9b-eksg} & map & 7{,}317 & 0 & 2023-10 &  &  & n\\
Historic Special Use Districts - 2010 & \texttt{kmx9-ph84} & map & 79 & 0 & 2023-10 &  &  & n\\
Historic Special Use Districts - 2011 & \texttt{st6c-ij8w} & map & 106 & 0 & 2023-10 &  &  & n\\
Historic Special Use Districts - 2012 & \texttt{9d9m-eyy3} & map & 148 & 0 & 2023-10 &  &  & n\\
Historic Special Use Districts - 2015 & \texttt{et47-hd49} & map & 146 & 0 & 2023-10 &  &  & n\\
Historic Zoning Districts - 1998 & \texttt{aypm-4d84} & map & 152{,}569 & 0 & 2023-10 &  &  & n\\
Historic Zoning Districts - 2000 & \texttt{aksh-67x3} & map & 153{,}292 & 0 & 2023-10 &  &  & n\\
Historic Zoning Districts - 2001 & \texttt{pdvd-w2q4} & map & 153{,}669 & 0 & 2023-10 &  &  & n\\
Historic Zoning Districts - 2002 & \texttt{ftvx-vtyc} & map & 153{,}663 & 0 & 2023-10 &  &  & n\\
Historic Zoning Districts - 2003 & \texttt{kspe-fmej} & map & 153{,}716 & 0 & 2023-10 &  &  & n\\
Historic Zoning Districts - 2004 & \texttt{f883-k722} & map & 191{,}062 & 0 & 2023-10 &  &  & n\\
Historic Zoning Districts - 2005 & \texttt{u8bf-s4hg} & map & 196{,}344 & 0 & 2023-10 &  &  & n\\
Historic Zoning Districts - 2007 & \texttt{rxim-5dpc} & map & 153{,}822 & 0 & 2023-10 &  &  & n\\
Historic Zoning Districts - 2009 & \texttt{u4br-5hb8} & map & 23{,}355 & 0 & 2023-10 &  &  & n\\
Historic Zoning Districts - 2010 & \texttt{w3j2-4hed} & map & 9{,}897 & 0 & 2023-10 &  &  & n\\
Historic Zoning Districts - 2012 & \texttt{9xqx-n98b} & map & 9{,}889 & 0 & 2023-10 &  &  & n\\
Historic Zoning Districts - 2013 & \texttt{manb-7mxp} & map & 9{,}982 & 0 & 2023-10 &  &  & n\\
Historic Zoning Districts - 2014 & \texttt{jn35-yfmd} & map & 9{,}981 & 0 & 2023-10 &  &  & n\\
Historic Zoning Districts - Jan 2015 & \texttt{utj8-tgqr} & map & 9{,}978 & 0 & 2023-10 &  &  & n\\
Historic Zoning Districts - March 2005 & \texttt{d7vm-pqzv} & map & 196{,}344 & 0 & 2023-10 &  &  & n\\
Historical Zoning Districts - April 2000 & \texttt{prs8-k8k3} & map & 153{,}292 & 0 & 2023-10 &  &  & n\\
Landmark Districts & \texttt{vnrd-fpg7} & map & 20 & 0 & 2023-10 &  &  & n\\
Maher & \texttt{hqsk-4xmh} & map & 974 & 0 & 2026-09 &  &  & n\\
Map of 2020 Census Tracts Assigned to Analysis Neighborhoods & \texttt{rqw6-h7c5} & map & 242 & 0 & 2022-07 &  &  & n\\
Map of Historic Districts & \texttt{y75h-nbt2} & map & 204 & 0 & 2026-09 &  &  & n\\
Map of Historic Resource Status by Property & \texttt{y8v4-jjwg} & map & 163{,}278 & 0 & 2026-09 &  &  & n\\
Map of Historic Special Use Districts - 2014 & \texttt{uihu-8bih} & map & 150 & 0 & 2023-10 &  &  & n\\
Map of Historic Special Use Districts -2013 & \texttt{r6jz-e8rw} & map & 147 & 0 & 2023-10 &  &  & n\\
Map of Historic Zoning Districts - 2006 & \texttt{v9w2-c8v7} & map & 23{,}132 & 0 & 2023-10 &  &  & n\\
Map of Historic Zoning Districts - 2008 & \texttt{ahwu-55yv} & map & 153{,}355 & 0 & 2016-08 &  &  & n\\
Map of Historic Zoning Districts - 2011 & \texttt{by3b-5pje} & map & 9{,}880 & 0 & 2023-10 &  &  & n\\
Map of Landmarks Listed in Article 10 of the San Francisco Planning Code & \texttt{hycf-nc3x} & map & 370 & 0 & 2026-09 &  &  & n\\
Map of Privately Owned Public Open Spaces & \texttt{33jz-gqck} & map & 81 & 0 & 2026-09 &  &  & n\\
Map of SF Urban Tree Canopy & \texttt{55pv-5zcc} & map & 289{,}219 & 0 & 2016-07 &  &  & n\\
Map of San Francisco Land Use - 2023 & \texttt{k8rg-ihdq} & map & 154{,}476 & 0 & 2024-03 &  &  & n\\
Map of slopes of 20\% or greater & \texttt{3vv2-nvev} & map & 28{,}055 & 0 & 2017-08 &  &  & n\\
Neighborhood-Specific Impact Fee Areas & \texttt{4x4u-9ser} & map & 33 & 0 & 2024-11 &  &  & n\\
Planning Department Neighborhood Quadrants & \texttt{x89n-qawp} & map & 4 & 0 & 2016-03 &  &  & n\\
Planning Districts & \texttt{ttns-6zj3} & map & 17 & 0 & 2023-10 &  &  & n\\
Planning Neighborhood Groups Map & \texttt{iacs-ws63} & map & 37 & 0 & 2023-10 &  &  & n\\
SF Find Neighborhoods & \texttt{pty2-tcw4} & map & 117 & 0 & 2023-10 &  &  & n\\
San Francisco Development Pipeline & \texttt{k55i-dnjd} & map & 1{,}917 & 0 & 2026-08 &  &  & n\\
San Francisco Environmental Justice Communities Map & \texttt{v25p-yprq} & map & 2{,}702 & 0 & 2023-07 &  &  & n\\
Zoning Map - Height and Bulk Districts & \texttt{kxax-c386} & map & 5{,}292 & 0 & 2026-09 &  &  & n\\
Zoning Map - Special Sign Districts - Scenic Streets & \texttt{4hx7-bwhg} & map & 12 & 0 & 2024-05 &  &  & n\\
{[}ARCHIVED{]} Map of Land Use - 2020 & \texttt{us3s-fp9q} & map & 155{,}395 & 0 & 2023-10 &  &  & n\\
{[}DEPRECATED{]} Historic Resources & \texttt{jfw5-qwcg} & map & 0 & 0 & 2026-05 &  &  & n\\
\multicolumn{9}{@{}l}{\textbf{Municipal Transportation Agency} (15)}\\
Equity Strategy Neighborhoods & \texttt{f8bh-m8ez} & dataset & 9 & 11 & 2024-07 & polygon & date & n\\
MTA.\allowbreak{}equitystrategy\_\allowbreak{}nhoods\_\allowbreak{}data & \texttt{qqn9-j3rz} & dataset & 124 & 15 & 2018-10 & point & date & n\\
Mid-Block Traffic Calming Areas & \texttt{abhw-ffzx} & dataset & 1{,}686 & 23 & 2026-05 & line & date & n\\
On-street Carshare Parking & \texttt{g2t6-cyw6} & dataset & 175 & 44 & 2024-01 & point & date & n\\
SFMTA - Mobility Data Specification & \texttt{h5q5-vfti} & dataset & 63 & 19 & 2025-09 & polygon & date & n\\
Sidewalk Widths (2014) & \texttt{4g86-grxu} & dataset & 16{,}174 & 14 & 2024-06 & line & date & n\\
Slow Streets & \texttt{hkz3-itiu} & dataset & 240 & 23 & 2024-01 & line & date & n\\
Traffic Signals & \texttt{ybh5-27n2} & dataset & 1{,}507 & 55 & 2024-05 & point & date & n\\
Travel Decision Survey 2021 & \texttt{b7eg-vqw2} & dataset & 721 & 263 & 2022-04 &  & date & n\\
SFMTA General Transit Feed Specification (GTFS) Production & \texttt{dni7-qpv3} & file & --- & 0 & 2022-06 &  &  & n\\
Travel Decision Survey Data 2013 & \texttt{j6i5-prai} & file & --- & 0 & 2015-09 &  &  & n\\
Equity Strategy Neighborhoods & \texttt{u7ag-wi48} & map & 9 & 0 & 2024-07 &  &  & n\\
MTA.\allowbreak{}paintedsafetyzones & \texttt{qqnk-bwbc} & map & 124 & 0 & 2018-10 &  &  & n\\
Map of Mid-Block Traffic Calming Areas & \texttt{5wa2-2y3e} & map & 1{,}686 & 0 & 2026-05 &  &  & n\\
Map of On-street Carshare Parking & \texttt{wr3x-ijzr} & map & 175 & 0 & 2024-01 &  &  & n\\
\multicolumn{9}{@{}l}{\textbf{Building Inspection} (14)}\\
Permits Issued (Jan 2013 - Current) & \texttt{n644-pp3v} & chart & --- & 0 & 2026-09 &  &  & n\\
Dwelling Unit Completion Counts by Building Permit$^{\dagger}$ & \texttt{j67f-aayr} & dataset & 2{,}541 & 5 & 2026-09 &  & permit, date & p\\
Building Permit Addenda with Routing & \texttt{87xy-gk8d} & filtered view & 3{,}998{,}159 & 20 & 2026-09 &  & permit, date & c\\
Building Permit Application Issuance Metrics & \texttt{gzxm-jz5j} & filtered view & 148{,}951 & 24 & 2026-09 &  & block+lot, permit, date & p\\
Building Permit Application Review Metrics & \texttt{5bat-azvb} & filtered view & 453{,}703 & 33 & 2026-09 &  & block+lot, permit, date & p\\
Building Permits$^{\dagger}$ & \texttt{i98e-djp9} & filtered view & 1{,}295{,}168 & 53 & 2026-09 & point & block+lot, case/\allowbreak record, permit, date & c\\
Building Permits Deduplicated on Primary Address & \texttt{f2jc-ivnc} & filtered view & 1{,}149{,}358 & 51 & 2026-09 & point & block+lot, case/\allowbreak record, permit, date & p\\
Building Permits before January 1, 2013 & \texttt{4jpb-z4kk} & filtered view & 849{,}624 & 41 & 2026-09 &  & block+lot, case/\allowbreak record, permit, date & n\\
Building Permits filed on or after January 1, 2013 & \texttt{p4e4-a5a7} & filtered view & 445{,}544 & 51 & 2026-09 & point & block+lot, case/\allowbreak record, permit, date & n\\
Building Permits with Contractor, Subcontractor, Architect, Designer, or Engineer & \texttt{cw8k-gwb7} & filtered view & 853{,}025 & 30 & 2026-09 & point & permit, date & p\\
Building Permits with Permit Contacts & \texttt{9itm-3rmi} & filtered view & 1{,}725{,}431 & 74 & 2026-09 & point & block+lot, case/\allowbreak record, permit, date & p\\
Map of Building Permits & \texttt{jgws-wax9} & map & 1{,}295{,}168 & 0 & 2026-09 &  &  & n\\
Map of Issued and Reinstated Building Permits & \texttt{4yy5-v29h} & map & 1{,}295{,}168 & 0 & 2026-09 &  &  & n\\
Building Permits Lookup by Contractor, Subcontractor, Architect, Designer, or Engineer & \texttt{kvek-u79k} & story & --- & 0 & 2025-11 &  &  & n\\
\multicolumn{9}{@{}l}{\textbf{Office of Resilience and Capital Planning} (14)}\\
100-Year Storm + 24" Sea Level Rise & \texttt{esku-ejgv} & dataset & 7 & 5 & 2020-02 & polygon &  & p\\
100-Year Storm Flood Risk Zone (July 2022) & \texttt{jzu3-4yxp} & dataset & 1 & 5 & 2022-08 & polygon & date & p\\
100-year storm + 66 inches Sea Level Rise & \texttt{6ann-usi8} & dataset & 1 & 4 & 2020-02 & polygon &  & p\\
Dam or Reservoir Failure Hazard Zone & \texttt{7crf-46s6} & dataset & 7 & 11 & 2020-02 & polygon & date & p\\
FEMA FIRM Flood Hazards (Coastal) - 2021 Update & \texttt{jyce-e25k} & dataset & 222 & 22 & 2021-01 & polygon &  & p\\
HAQR Priority Green Infrastructure Zones & \texttt{nn26-kuy2} & dataset & 1 & 2 & 2024-06 & polygon &  & p\\
Hayward M7.\allowbreak{}0 Earthquake Scenario & \texttt{d6pi-gknj} & dataset & 26 & 16 & 2020-02 & polygon & date & p\\
Landslide Susceptibility Hazard Zones & \texttt{bna4-itif} & dataset & 8 & 11 & 2020-02 & polygon & date & p\\
San Andreas M7.\allowbreak{}8 Earthquake Scenario & \texttt{akmp-ztcg} & dataset & 27 & 16 & 2020-02 & polygon & date & p\\
Soil Liquefaction Hazard Zone & \texttt{i4t7-35u3} & dataset & 2 & 4 & 2020-02 & polygon &  & p\\
Tsunami Inundation Hazard Zone (2021 Update) & \texttt{7p2k-c3z6} & dataset & 1 & 11 & 2021-08 & polygon &  & p\\
Tsunami Inundation Hazard Zone (Outdated, 2015) & \texttt{q7uq-hghg} & dataset & 3 & 12 & 2020-02 & polygon & date & n\\
Capital Budget for General Fund Departments (2005-06 to present) & \texttt{qf4b-pv4d} & filtered view & 3{,}754 & 22 & 2026-09 &  & date & n\\
Map of Tsunami Inundation Hazard Zone (2021 Update) & \texttt{xidx-bh43} & map & 1 & 0 & 2021-08 &  &  & n\\
\multicolumn{9}{@{}l}{\textbf{Public Works} (11)}\\
Covid-19 Pit Stops & \texttt{mr6h-cr3u} & dataset & 16 & 13 & 2026-09 & point & date & n\\
Curb Ramps & \texttt{ch9w-7kih} & dataset & 50{,}098 & 40 & 2026-04 & point & date & n\\
Large Utility Excavation Permits & \texttt{i926-ujnc} & dataset & 3{,}036 & 17 & 2026-09 & point & permit, date & n\\
Right of Way Exception Data & \texttt{yrgu-vakm} & dataset & 106{,}191 & 14 & 2026-09 &  & block+lot, date & n\\
SF County Line 2020 Reestablishment & \texttt{hmft-3hxr} & dataset & 1 & 5 & 2026-03 & line & date & n\\
Street Tree Removal Notifications & \texttt{qrwx-q4gg} & dataset & 5{,}701 & 25 & 2026-09 & point & date & n\\
Street-Use Permits & \texttt{b6tj-gt35} & dataset & 1{,}228{,}392 & 31 & 2026-09 & point & permit, date & n\\
Utility Excavation Moratorium Streets & \texttt{5wbp-dwzt} & dataset & 3{,}390 & 19 & 2026-09 & point & date & n\\
Utility Excavation Permits & \texttt{smdf-6c45} & dataset & 4{,}057 & 11 & 2026-09 &  & permit, date & n\\
Map of Curb Ramps & \texttt{2y4g-3ig2} & map & 50{,}098 & 0 & 2026-04 &  &  & n\\
Map of SF County Line 2020 & \texttt{jxv4-h25i} & map & 1 & 0 & 2026-03 &  &  & n\\
\multicolumn{9}{@{}l}{\textbf{GSA - City Administrator's Office} (8)}\\
108" Inundation Vulnerability Zone Line (Sea Level Rise + 100YR Flood Event) & \texttt{92e4-7ptg} & dataset & 1{,}005 & 15 & 2017-03 & line &  & p\\
2015 Project Appendix & \texttt{xpdk-pzjd} & dataset & 126 & 33 & 2015-08 & point &  & n\\
City Construction Projects to Proceed During Shelter-in-place & \texttt{x83f-ux29} & dataset & 131 & 4 & 2020-04 &  &  & n\\
Parcels with overlay attributes & \texttt{9grn-xjpx} & dataset & 236{,}556 & 17 & 2026-09 & point & blklot, block+lot, date & p\\
Tall Building Inventory & \texttt{5kya-mfst} & dataset & 180 & 45 & 2019-07 & polygon & date & n\\
City Facilities - Planning Jurisdiction or Leased & \texttt{nfxf-sv4m} & filtered view & 2 & 14 & 2026-09 &  & blklot & n\\
2015 Project Appendix Map & \texttt{dyxq-hen7} & map & 126 & 0 & 2015-08 &  &  & n\\
Map of 108" Inundation Vulnerability Zone Line (Sea Level Rise + 100YR Flood Event) & \texttt{ff7h-99em} & map & 1{,}005 & 0 & 2017-03 &  &  & n\\
\multicolumn{9}{@{}l}{\textbf{Environment} (7)}\\
Existing Buildings Energy Performance Ordinance Report & \texttt{96ck-qcfe} & dataset & 30{,}544 & 34 & 2026-09 & point & blklot, date & n\\
San Francisco Communitywide Greenhouse Gas Inventory & \texttt{btm4-e4ak} & dataset & 834 & 11 & 2024-06 &  &  & n\\
San Francisco Municipal Natural Gas Equipment Inventory & \texttt{vc6r-v7av} & dataset & 1{,}612 & 23 & 2026-09 & point & blklot, date & n\\
San Francisco Plant Finder Data & \texttt{vmnk-skih} & dataset & 629 & 43 & 2025-07 &  &  & n\\
Existing Buildings - Basic Info and Audit Compliance Status & \texttt{vgqy-2ca4} & filtered view & 2{,}494 & 17 & 2026-09 & point & blklot, date & n\\
Existing Buildings - Benchmark Reports & \texttt{4ua7-5sfx} & filtered view & 30{,}544 & 32 & 2026-09 & point & blklot, date & n\\
Map of natural gas-fueled equipment in municipally owned buildings & \texttt{h5j5-ymyb} & map & 1{,}612 & 0 & 2026-09 &  &  & n\\
\multicolumn{9}{@{}l}{\textbf{(not stated)} (6)}\\
Zoning Map - Special Use Districts$^{\dagger}$ & \texttt{rxfc-8aap} & dataset & 125 & 7 & 2024-05 & polygon & date & c\\
Zoning Map - Zoning Districts$^{\dagger}$ & \texttt{xzez-p3nc} & dataset & 1{,}647 & 12 & 2024-02 & polygon & date & c\\
Article 11 Rating & \texttt{ee6v-jtbt} & map & 786 & 0 & 2024-02 &  &  & n\\
Zoning Map - Special Use Districts & \texttt{ry69-hnut} & map & 125 & 0 & 2024-05 &  &  & n\\
Zoning Map - Zoning Districts & \texttt{ibu8-4ccn} & map & 1{,}647 & 0 & 2024-02 &  &  & n\\
Dataset Inventory Explorer & \texttt{dztp-vt7z} & story & --- & 0 & 2025-07 &  &  & n\\
\multicolumn{9}{@{}l}{\textbf{Controller} (5)}\\
Budget & \texttt{xdgd-c79v} & dataset & 408{,}232 & 24 & 2026-08 &  & date & n\\
Child Abuse Prevention Services in San Francisco & \texttt{3had-h899} & dataset & 375 & 61 & 2019-12 &  &  & n\\
Employee Compensation & \texttt{88g8-5mnd} & dataset & 1{,}180{,}053 & 26 & 2026-09 &  & date & n\\
General Obligation (GO) Bond Projects & \texttt{d3dc-v5yr} & dataset & 305 & 11 & 2026-04 &  & date & n\\
San Francisco Citywide Performance Metrics & \texttt{6dpt-7w23} & dataset & 10{,}164 & 15 & 2026-06 &  & date & n\\
\multicolumn{9}{@{}l}{\textbf{Technology} (5)}\\
City Lands & \texttt{gtnh-hgvs} & dataset & 1{,}035 & 12 & 2026-09 & polygon & date & p\\
Parcels – Active and Retired$^{\dagger}$ & \texttt{acdm-wktn} & dataset & 236{,}556 & 37 & 2026-09 & polygon & blklot, block+lot, date & c\\
San Francisco ZIP Codes & \texttt{uq3t-6t53} & dataset & 32 & 13 & 2023-10 & polygon &  & p\\
Parcels - Active and Retired (no geometry) & \texttt{8jwb-2stv} & filtered view & 236{,}556 & 36 & 2026-09 &  & blklot, block+lot, date & n\\
San Francisco ZIP Codes & \texttt{srq6-hmpi} & map & 32 & 0 & 2023-10 &  &  & n\\
\multicolumn{9}{@{}l}{\textbf{Human Resources} (4)}\\
City Employment Classification Definitions & \texttt{7br7-cj6v} & dataset & 1{,}217 & 16 & 2022-12 &  & date & n\\
City Employment Compensation Plan & \texttt{h3jp-c25d} & dataset & 20{,}637 & 32 & 2026-06 &  & date & n\\
Management Classification and  Compensation Plan Pay Ranges & \texttt{mx8n-vbe2} & dataset & 1{,}510 & 16 & 2024-10 &  & date & n\\
Permanent Civil Service (PCS) Appointments & \texttt{qfj2-m7ux} & dataset & 26{,}131 & 9 & 2026-09 &  & date & n\\
\multicolumn{9}{@{}l}{\textbf{Ethics Commission} (4)}\\
SFEC Form 3410A - Registration Report for Permit Consultants & \texttt{umwe-sn9p} & dataset & 170 & 14 & 2026-08 &  & date & p\\
SFEC Form 3410B - Quarterly Report for Permit Consultants Filings & \texttt{k2rm-926a} & dataset & 1{,}531 & 18 & 2026-07 &  & date & p\\
SFEC Form 3500 - Disclosure Report for Developers of Major City Projects Filings & \texttt{8je3-gv3q} & dataset & 174 & 27 & 2026-07 &  & case/\allowbreak record, date & p\\
SFEC Form 3500 - Disclosure Report for Developers & \texttt{bk74-rima} & filtered view & 174 & 26 & 2026-07 &  & case/\allowbreak record, date & p\\
\multicolumn{9}{@{}l}{\textbf{Mayor} (4)}\\
Mayor's Office of Housing and Community Development Affordable Housing Pipeline$^{\dagger}$ & \texttt{aaxw-2cb8} & dataset & 190 & 68 & 2026-05 & point & block+lot, case/\allowbreak record, permit, date & p\\
Mayor's Office of Housing and Community Development Affordable Housing Portfolio & \texttt{pyxv-n29e} & dataset & 850 & 34 & 2026-05 & point & date & p\\
Mayor's Office of Housing and Community Development Non-Unit Inclusionary Housing Alternatives & \texttt{g6kh-9pnv} & dataset & 153 & 29 & 2026-05 &  & case/\allowbreak record, permit, date & p\\
Map of Mayor's Office of Housing and Community Development Affordable Housing Portfolio & \texttt{k8na-jk8h} & map & 850 & 0 & 2026-05 &  &  & n\\
\multicolumn{9}{@{}l}{\textbf{Health Service System} (3)}\\
HSS View-Medical Enrollment by Employees & \texttt{v8fd-x8ui} & filtered view & 443 & 8 & 2026-03 &  &  & n\\
HSS-View-Enrollment\_\allowbreak{}all\_\allowbreak{}Lives\_\allowbreak{}5y & \texttt{phns-r8sy} & filtered view & 70 & 7 & 2026-03 &  &  & n\\
View of Medical and Dental Enrollment by Union for City \& County of SF Employees(CCSF) & \texttt{42ay-783m} & filtered view & 1{,}490 & 8 & 2024-06 &  &  & n\\
\multicolumn{9}{@{}l}{\textbf{DataSF} (3)}\\
Dataset inventory & \texttt{y8fp-fbf5} & dataset & 1{,}336 & 23 & 2026-09 &  & date & p\\
Datasets in 'Scoping' or 'In Progress' phase & \texttt{xu9j-dnfu} & filtered view & 64 & 10 & 2026-09 &  & date & p\\
Land Use with Active Parcel Information & \texttt{h5vv-2u24} & filtered view & 224{,}322 & 23 & 2024-03 & polygon & blklot & p\\
\multicolumn{9}{@{}l}{\textbf{Permit Center} (2)}\\
Permit\allowbreak{}SF Permitting Data & \texttt{tyz3-vt28} & dataset & 4{,}999 & 24 & 2026-09 &  & blklot, case/\allowbreak record, date & p\\
Building Permit Completeness Check Review Metrics & \texttt{abh5-gwaq} & filtered view & 6{,}618 & 10 & 2026-09 &  & date & p\\
\multicolumn{9}{@{}l}{\textbf{Public Health} (2)}\\
Maternal, Child, and Adolescent Health Needs Assessment, 2023-2024 & \texttt{iqtk-etij} & dataset & 62{,}714 & 28 & 2025-08 &  &  & n\\
San Francisco Flood Health Vulnerability & \texttt{cne3-h93g} & dataset & 578 & 27 & 2016-03 &  &  & n\\
\multicolumn{9}{@{}l}{\textbf{Emergency Management} (2)}\\
COVID-19 Tenderloin Plan Zones & \texttt{xw26-dnmx} & dataset & 6 & 4 & 2020-06 & polygon &  & n\\
Tenderloin COVID-19 Response: Sanitation Map & \texttt{fi63-bg5f} & map & 6 & 0 & 2020-06 &  &  & n\\
\multicolumn{9}{@{}l}{\textbf{Police Department} (2)}\\
Police Department Incident Reports: Historical 2003 to May 2018 & \texttt{tmnf-yvry} & dataset & 2{,}071{,}736 & 36 & 2025-06 & point & date & n\\
SF Crime Heat Map & \texttt{q6gg-sa2p} & map & 2{,}071{,}736 & 0 & 2025-06 &  &  & n\\
\multicolumn{9}{@{}l}{\textbf{Economic \& Workforce Development} (1)}\\
Invest in Neighborhood Areas & \texttt{667d-aah6} & dataset & 25 & 11 & 2016-06 & polygon &  & n\\
\multicolumn{9}{@{}l}{\textbf{Retirement System} (1)}\\
Retiree Pensions Annual Benefit Received & \texttt{84zg-myyk} & dataset & 28{,}993 & 5 & 2023-11 &  & date & n\\
\multicolumn{9}{@{}l}{\textbf{Digital Services} (1)}\\
Inventory of web domains managed by City and County departments & \texttt{m26n-jhf2} & dataset & 432 & 8 & 2026-09 &  & date & n\\
\multicolumn{9}{@{}l}{\textbf{Other} (1)}\\
ARCHIVED: Tenderloin Tent Count & \texttt{ubjr-y9qw} & dataset & 339 & 2 & 2023-07 &  & date & n\\
\bottomrule\end{longtable}}

\begin{table}[htbp]\centering
\caption{Negative results. A row here is a search that was run and is not to be run again.}\label{tab:negatives}
\begin{tabular}{@{}L{3.6cm}L{4.6cm}L{6.0cm}@{}}\toprule
Sought & Where & Result\\\midrule
A dataset of Commission actions, motions or conditions & all 273 catalogue results; Planning's 109 entries in the dataset inventory, published or not & none; the inventory entry that use the words: Privately Owned Public Open Spaces\\[2pt]
The catalogue by the old host name & \texttt{search\_context} or \texttt{domains} = \texttt{data.sfgov.org} & 273 with \texttt{search\_context}, the same assets as the new host; 2 with \texttt{domains} alone, as for the new host --- not zero\\[2pt]
Every result, by paging the default order & the Discovery API, relevance order & 273 rows holding 261 distinct assets; page in \texttt{order=dataset\_id}\\[2pt]
Planning permitting data & \texttt{kncr-c6jw} [Deprecated] & 0 rows; its description names the two records tables, both cached, as its replacement\\[2pt]
Historic resources & \texttt{njrt-gtwr} [Deprecated] & 0 rows; its description names \texttt{3tsw-4idn} as its replacement\\[2pt]
Extra zoning vintages & \texttt{5uzm-v8n8} (March 2005), \texttt{cfhj-8c6b} (April 2000) & 196{,}344 and 153{,}292 rows, the row counts of the cached 2005 and 2000 vintages (196{,}344, 153{,}292); probably duplicates; compared on row count only\\[2pt]
A published history of the Development Pipeline & the dataset inventory & 4 quarterly snapshots are listed as held and not published\\[2pt]
\bottomrule\end{tabular}\end{table}

\begin{longtable}{@{}L{1.6cm}L{1.9cm}L{4.2cm}L{6.6cm}@{}}
\caption{What the Planning Department lists in the city's dataset inventory (\texttt{y8fp-fbf5}) and has not published. Descriptions are the inventory's own, cut to their first sentence or so.}\label{tab:inventory}\\\toprule
Inventory & Status & Dataset & What the inventory says\\\midrule\endfirsthead
\toprule Inventory & Status & Dataset & What the inventory says\\\midrule\endhead
\texttt{CPC-0020} & In Progress & Parcels with Active Planning Complaints (for OpenGov) & nan\\
\texttt{CPC-0021} & In Progress & Addresses in SF Port Jurisdiction (for OpenGov) & nan\\
\texttt{CPC-0182} & In Progress & San Francisco Environmental Justice Communities Map & The Environmental Justice Communities Map (“EJ Communities Map”) describes areas of San Francisco that have higher pollution and are predominately low-income. This map is\dots\\
\texttt{CPC-0092} & Internal & SF Development Pipeline 2019 Q4 & Snapshot of San Francisco Development Pipeline. Tracking of construction and entitlement activity based on data from Department of Building Inspection's Permit Tracking a\dots\\
\texttt{CPC-0130} & Internal & Urban Bird Refuge & The adopted Standards for Bird-Safe Buildings explains the documented risks that structures present to birds. Over thirty years of research has proven the risk to be biol\dots\\
\texttt{CPC-0162} & Internal & SF Development Pipeline 2022 Q1 {[}REVISED{]} & Note: This dataset (SF Development Pipeline 2022 Q1) was updated on Dec 21, 2022 to correct missing data from the previous version. A. SUMMARY Snapshot of San Francisco's\dots\\
\texttt{CPC-0163} & Internal & SF Development Pipeline 2022 Q3 & Snapshot of San Francisco's Development Pipeline. Tracking of construction and entitlement activity based on data from Department of Building Inspection's Permit Tracking\dots\\
\texttt{CPC-0165} & Internal & SF Development Pipeline 2022 Q2 & Snapshot of San Francisco's Development Pipeline. Tracking of construction and entitlement activity based on data from Department of Building Inspection's Permit Tracking\dots\\
\texttt{CPC-0169} & Internal & Priority Processing Programs & Data from the Priority Processing Programs table contained within the PRJ record type of the Planning Department's Accela permitting system. Types of Priority Processing \dots\\
\texttt{CPC-0170} & Internal & Housing Programs & Dataset of Housing Programs tracked by PRJ and PRL record types in the Planning Department's Accela permitting system. The Accessory Dwelling Unit (ADU) program is tracke\dots\\
\texttt{CPC-0175} & Internal & Planning Department Project Review Events & This dataset contains event-level data from the Planning Department's Accela software. The data is detailed at both the "project" and "event" levels. In some cases, metri\dots\\
\texttt{CPC-0176} & Internal & Planning Department Projects & This dataset contains project-leve data from the Planning Department's Accela software. In some cases, metrics related to the project timeline (such as duration) are prec\dots\\
\texttt{CPC-0179} & Internal & Historic Preservation Survey Units & This dataset is an export from the SF Planning Department's Arches implementation (https://sfculturalheritage.org/search) showing survey results of the San Francisco Cult\dots\\
\texttt{CPC-0180} & Internal & Planning Notices of Violation & This dataset is an extract from Accela that shows active notices of violation sent to the Planning Department. B. HOW THE DATASET IS CREATED This dataset is extracted fro\dots\\
\texttt{CPC-0019} & Not Published & Open Spaces - City Planning & Map of publicly accessible open spaces in San Francisco\\
\texttt{CPC-0022} & Not Published & Landmark Streets & Map of landmark streets. Landmarks and Landmark Districts are defined and listed in Article 10 of the San Francisco Planning Code.\\
\texttt{CPC-0127} & Not Published & National Register Districts & Districts that are either listed or determined eligible to be listed on the National Register by San Francisco Planning Department preservation staff. Data is in a zipped\dots\\
\texttt{CPC-0128} & Not Published & California Register Districts & Districts that are either listed or determined eligible to be listed on the California Register by San Francisco Planning Department preservation staff. Data is in a zipp\dots\\
\texttt{CPC-0157} & Not Published & Planning Record ID to APN Crosswalk & Reference dataset that has the APNs related to a planning project\\
\texttt{CPC-0178} & Not Published & Short Term Rental Permits & https://sfplanning.org/office-short-term-rentals\\
\bottomrule\end{longtable}


\paragraph{Two ways the catalogue gives a wrong answer without an error.}
\begin{itemize}
\item \textbf{Paging the default order drops assets.} Results come back in relevance order,
and a relevance-ranked list re-ranks between requests: paged to exhaustion that way, the
query returned \catRelPagingRows\ rows holding \catRelPagingDistinct\ distinct assets ---
some twice, some never. Ordered by identifier
(\texttt{order=dataset\_id}) it returns every asset once. Everything in this memo was paged
that way.
\item \textbf{Which parameter names the host matters more than which host is named.}
Table~\ref{tab:hosttest} runs the same query six ways. Passed as \texttt{search\_context},
the old host \texttt{data.sfgov.org} returns \catOldHostCount\ results --- \catOldHostSameIds\
assets as the new host. Passed as \texttt{domains} alone, the new host returns \catDomainsOnly\ and the old
\catDomainsOnlyOld. The data-acquisition memo's finding that ``a catalogue search against the
old host returns zero hits'' did not reproduce on \catRetrieved; whether the catalogue has
changed since that memo's run or the earlier test passed the host differently is not known,
and that memo now carries a dated note saying so.
\end{itemize}

\section{What each asset carries}

For each of the \catTotal\ assets the catalogue record and the view's own metadata give its
identifier, name, type, parent dataset where it is a derived view, publishing department,
category, tags, creation and last-update dates, stated update frequency, column names and
types, and the first part of its description; \texttt{SELECT count(*)} against the resource
endpoint gives its row count. \catWithRows\ assets have rows of their own, together
\catRowsTotal\ of them; the rest are links, stories, files and a chart. \catGeom\ carry a
geometry column. All of this is in \texttt{catalogue\_planning.csv} (\S\ref{sec:repro}).

\paragraph{Join keys, by name.} A column name is a claim about the rows, not proof of what
they hold, and the table marks keys ``by name'' for that reason. \catKeyAny\ assets carry at
least one of the pipeline's keys: \catKeyBlklot\ a \texttt{blklot}-style parcel identifier,
\catKeyBlockLot\ separate block and lot columns, \catKeyCase\ a planning case or record
number, \catKeyPermit\ a permit number. The name-matching has one known false friend: DBI's
\texttt{record\_id} is DBI's own row identifier, not a Planning record, and it makes every
view of DBI's permits look as though it carries a case number.

\paragraph{Departments and categories.} Table~\ref{tab:depts} gives the \catDepartments\
publishers. The Planning Department publishes \catPlanningDept\ of the \catTotal; the rest
matched on the word ``planning'' somewhere in their text --- capital planning, transport
planning, workforce planning --- and include salary schedules, pension payments and police
incident reports. Table~\ref{tab:cats} gives the portal's own categories, which are
unassigned on a large share of assets and are not a useful filter.

\paragraph{Staleness.} \catStale\ assets have not been updated in the \catStaleYears\ years
before retrieval (Table~\ref{tab:stale}). For most of them that is by design --- a 2009
zoning vintage is not meant to change --- and staleness is a problem only for the assets
that claim to be current.

\section{Relevance, and what is already cached}

Each asset carries a flag --- \emph{core}, \emph{possibly useful} or \emph{not relevant} ---
and a one-sentence reason. \textbf{The flags are judgement, not measurement}: they were
written by hand, after reading every asset's name, description and columns, into
\texttt{relevance.csv} beside this memo, and the script only reads them. \catCore\ assets are
flagged core, \catPossible\ possibly useful and \catNotRelevant\ not relevant. \catCached\ of
the \catTotal\ are already cached by \texttt{acquire\_external\_data.py} or
\texttt{analyze\_permits.py}, and are marked with a dagger in Table~\ref{tab:full}; the check
is against the identifiers those scripts fetch, which include every source in Table~1 of the
data-acquisition memo that the catalogue returns for this query.

Table~\ref{tab:possible} lists the \catPossibleNotCached\ assets flagged core or possibly
useful that are not cached. In rough order of what they would add:

\begin{enumerate}
\item \textbf{The delay the paper measures, from the city's own clocks.} The Planning
Department's \emph{Project Application Review metrics} (\texttt{d4jk-jw33}) give review-time
events per application, keyed on \texttt{b1\_alt\_id}, which is the Accela record number and
joins to the item table's case number the way the records tables do. DBI's \emph{Building
Permit Addenda with Routing} (\texttt{87xy-gk8d}) dates every review station a permit passes
through, Planning's included. Neither has been measured against anything in this pipeline,
and both bear directly on the delay memo and on the permit clock.
\item \textbf{Newer publications of cached layers.} A \emph{Zoning Map --- Zoning Districts}
(\texttt{3i4a-hu95}) updated at retrieval sits beside the cached \texttt{xzez-p3nc}, which
is in the stale list of Table~\ref{tab:stale}; the same is true of the special-use
districts. The zoning
panel carries the current layer forward from 2016, so the layer it carries should be checked
against the newer one before the panel is used for recent years.
\item \textbf{Vintages the panel lacks.} The dated special-use district layers of 2010--2015,
archived land-use vintages, and a set of scanned zoning maps from 1921 onward --- the only
pre-1998 zoning the portal holds, and only as images.
\item \textbf{Policy geographies and review triggers by parcel}: Affordable Housing Bonus
eligibility, Housing Sustainability Districts, historic status, the Maher contamination
sites, steep slopes and the air-quality exposure zone.
\item \textbf{The political economy of permits}: developer disclosures for major projects
keyed on the planning case number, registered permit consultants and their filings, and the
architects and expediters on each permit.
\end{enumerate}

\paragraph{What the city holds and has not published.} The city's dataset inventory lists
what departments hold whether or not it is on the portal. The Planning Department has
\catInvPlanning\ entries in it, \catInvPublished\ published and \catInvUnpublished\ not;
Table~\ref{tab:inventory} lists the unpublished ones. Worth a records request among them:
an event-level extract of the department's project review (the inventory's
\emph{Planning Department Project Review Events}), a record-to-parcel crosswalk, and
\catInvPipelineInternal\ quarterly snapshots of the Development Pipeline held internally ---
which is the pipeline \emph{panel} the data-acquisition memo reported as unavailable because
the published pipeline is a current snapshot.

\section{Negatives}

Table~\ref{tab:negatives} records what looked promising and was not, so the searches are
not repeated. The first row is the one the project cares about: \textbf{no asset in the
catalogue and no entry in the Planning Department's inventory is a record of Commission
actions, motions or conditions}. The table names the inventory entries that use those words
(\catInvCommission); read in full, none is a record of what the Commission decided.

\section{Established, and not established}

\textbf{Established.}
\begin{enumerate}
\item The catalogue returns \catTotal\ assets for ``planning'' on \catRetrieved; the web
figure of 155 matches the API's \catDatasets\ datasets, and the remainder is reconciled by
type in Table~\ref{tab:types}.
\item Relevance-ordered paging is not exhaustive; identifier-ordered paging is.
\item The old host name works as \texttt{search\_context}; \texttt{domains} alone does not
work for either host.
\item No Commission-actions dataset exists on the portal or in Planning's inventory.
\item \catPossibleNotCached\ uncached assets are judged core or possibly useful, led by two
review-time datasets and newer copies of two cached zoning layers.
\end{enumerate}

\textbf{Not established.}
\begin{enumerate}
\item Whether the web UI's 155 is a datasets count. It matches, and the page could not be
rendered here to confirm it.
\item Whether the old-host result changed since the data-acquisition memo's run.
\item What any uncached asset actually joins to. The keys are read from column names; no
join rate is measured here, and by this project's rule nothing in
Table~\ref{tab:possible} is ``found'' until one is.
\item Whether the two zoning layers differ in substance. They differ in size and date; the
polygons have not been compared.
\item The relevance flags, which are one reader's judgement.
\end{enumerate}

\section{Reproducing this}
\label{sec:repro}

\begin{verbatim}
python analyze_datasf_catalogue.py fetch    # catalogue, host test, inventory,
                                            # one view + count(*) per asset
python analyze_datasf_catalogue.py fetch --refresh   # re-query the catalogue
python analyze_datasf_catalogue.py report            # CSV, tables, macros
\end{verbatim}

\texttt{fetch} writes the raw catalogue answer, with its retrieval timestamp, the host test,
the per-type facet counts and a relevance-order paging pass, to
\texttt{\$MFHR\_DATA\_ROOT/}\allowbreak\texttt{external/datasf\_catalogue/catalogue\_planning\_<date>.json};
one enrichment file per asset to \texttt{enrich/}; and Planning's rows of the dataset
inventory to \texttt{inventory\_planning.json}. It pauses between requests to stay under two
a second on each host. \texttt{report} writes \texttt{catalogue\_planning.csv} (one row per
asset, every field above plus the flag and the columns the key match found) and the two
generated files this memo inputs. The relevance flags are the one hand-written input, in
\texttt{relevance.csv} beside this file; every number in the text is a macro.

\end{document}
